\documentclass[11pt]{article}

\usepackage[margin=1in]{geometry}
\usepackage{amsmath,amssymb,amsthm,mathtools,mathrsfs,bm}
\usepackage{enumitem}
\usepackage{booktabs}
\usepackage{microtype}
\usepackage{hyperref}

\allowdisplaybreaks

\newtheorem{theorem}{Theorem}[section]
\newtheorem{definition}[theorem]{Definition}
\newtheorem{proposition}[theorem]{Proposition}
\newtheorem{lemma}[theorem]{Lemma}
\newtheorem{corollary}[theorem]{Corollary}
\theoremstyle{remark}
\newtheorem{remark}[theorem]{Remark}
\newtheorem{example}[theorem]{Example}

\newcommand{\R}{\mathbb R}
\newcommand{\C}{\mathbb C}
\newcommand{\N}{\mathbb N}
\newcommand{\Z}{\mathbb Z}
\newcommand{\mcN}{\mathcal N}
\newcommand{\mcE}{\mathcal E}
\newcommand{\mcT}{\mathcal T}
\newcommand{\mcG}{\mathcal G}
\newcommand{\mcU}{\mathscr U}
\newcommand{\mcCaux}{\mathscr C}
\newcommand{\mcS}{\mathscr S}
\newcommand{\mcX}{\mathscr X}
\newcommand{\mcF}{\mathcal F}
\newcommand{\mcO}{\mathcal O}
\newcommand{\bbB}{\mathbb B}
\newcommand{\one}{\mathbf 1}
\newcommand{\tr}{\operatorname{tr}}
\newcommand{\Ree}{\operatorname{Re}}
\newcommand{\dd}{\,\mathrm d}
\newcommand{\rank}{\operatorname{rank}}
\newcommand{\diag}{\operatorname{diag}}
\newcommand{\Span}{\operatorname{span}}
\newcommand{\Ad}{\operatorname{Ad}}
\newcommand{\indeg}{\operatorname{indeg}}
\newcommand{\outdeg}{\operatorname{outdeg}}
\newcommand{\Vol}{\operatorname{Vol}}
\newcommand{\Id}{\mathrm{Id}}
\newcommand{\Ent}{\operatorname{Ent}}
\newcommand{\mcHh}{\mathfrak h_3}
\newcommand{\mcsu}{\mathfrak{su}(3)}

\title{From the Deterministic Regulator Data to a Gauge-Fixed Path Integral\\
for the Finite $SU(3)$ Yang--Mills Lattice Theory}
\author{}
\date{}

\begin{document}
\maketitle

\begin{abstract}
The deterministic regulator paper constructs, for each regulator size \(N\), the
gauge-geometric part of a finite regularized \(SU(3)\) data package
\[
\mathcal Q^{(N)}=
(\mcN^{(N)},\mcE^{(N)},\mathcal C_\square^{(N)},
\{U_3^{(N)}(e)\}_{e\in\mcE^{(N)}},
S_{\mathrm W}^{(N)})
\]
whose regulator geometry, sampled links, and Wilson action are fixed in
\cite{continuumqft}.  The paper completes the finite-dimensional data
package by adjoining the explicit fermionic and scalar actions written below and
performs the
quantization step: it replaces the sampled background links by independent
integration variables and builds the exact finite-dimensional Euclidean path
integral on that same regulator lattice.  We define the full configuration
space, the product Haar--Berezin--Lebesgue measure, the partition function, and
the vacuum expectation map.  We prove that the resulting finite-lattice theory
is well defined for the free massive scalar sector used in this paper.  We then
solve the gauge-redundancy problem exactly by
a rooted maximal-tree gauge: every gauge orbit meets the gauge slice once and
only once, the Faddeev--Popov determinant is identically $1$, and the ghost
sector is a decoupled finite Berezin integral.  The fermionic sector integrates
to an exact finite fermion determinant, and the scalar sector is shown to be
exactly Gaussian with an explicit lattice covariance.

The paper also establishes the Lie-algebraic structure of the discrete measure.
On each edge, Haar measure is smooth in exponential coordinates and agrees to
second order with Lebesgue measure on the Lie algebra.
The asymptotic regime relevant to this programme is the stationary
thermodynamic limit
\[
\pi_{N,1}\Rightarrow \pi_\infty,
\]
already proved in the companion convergence paper.  The convergence paper
\cite{convergence} proves this stationary mean-field limit, while
\cite{continuumqft} identifies the deterministic finite Wilson data and the
continuum functional \(S_{\mathrm{YM}}^{\mathrm{cont}}\) on
$\pi_\infty$.  The same convergence paper also proves the mean-field
logarithmic Sobolev inequality on $\pi_\infty$; by tensorization, the same
constant controls the finite product measures $\pi_\infty^{\otimes k}$ that
support the thermodynamic observables.  The paper supplies the exact
fixed-$N$ quantization, isolates a canonical continuum-indexed class of local
gauge-invariant plaquette observables together with its finite-regulator
approximation map, identifies the corresponding exact finite unnormalized
Euclidean functionals, and singles out a distinguished localized Wilson sector
whose sampled deterministic continuum limit is already identified in
\cite{continuumqft}.  The transferred equilibrium stiffness continues to
control the product measures supporting the thermodynamic plaquette
presentation.
\end{abstract}

\tableofcontents

\section{Introduction and motivation}
\label{sec:introduction}

\paragraph{From data to dynamics.}
The companion finite-regulator paper \emph{Local Continuum Yang--Mills Limit
from Particle Trajectories} constructs, for each regulator size, a finite
deterministic spacetime lattice
\[
\mathcal L=(\mcN,\mcE)
\]
whose site set and edge set form a deterministic coordinate graph, and whose
canonical plaquette family $\mathcal C_\square$ is the corresponding
plaquette set from \cite{continuumqft}.  On that lattice the same paper defines
sampled link transporters $U_3^{(N)}(e)\in SU(3)$ and the Wilson action
$S_{\mathrm W}$.  The paper places on the same regulator graph the
finite matter actions $S_{\mathrm F}$ and $S_\phi$ defined in
Definition~\ref{def:imported-package}.  Together these objects form the finite
data package $\mathcal Q^{(N)}$.  Up to that stage, however,
the link variables are deterministic functions of the continuum-seeded profile.
A quantum theory requires a different mathematical object: instead of one
distinguished configuration, it requires an ensemble weighted by an action and
summed over all admissible gauge fields and matter fields.  The paper
performs that passage from one explicit background configuration to a
finite-dimensional path integral on the same regulator lattice.

\paragraph{The quantization objective.}
The basic quantization step is conceptually simple and mathematically precise.
The background transporters $U_3^{(N)}(e)$ derived in \cite{continuumqft} are replaced
by independent variables
\[
U(e)\in SU(3),\qquad e\in\mcE,
\]
and the site colors $\mathbf c_x$ are retained, when desired, as auxiliary
fluctuating variables in the unit ball of $\C^3$.  Because the companion paper
derives no autonomous color-field energy beyond the induced links, the
auxiliary color sector will be normalized and dynamically inert in the
quantization performed here.  All dynamical content therefore sits in the gauge
links, fermions, and scalars.  The goal is to define a mathematically complete
finite partition function, remove gauge redundancy exactly, identify the
fermion determinant and scalar measure, and clarify in what sense the resulting
finite theory organizes the canonical local gauge-invariant observables used in
the thermodynamic continuum programme.

\begin{theorem}[Main finite-lattice quantization result]
\label{thm:main}
Fix the lattice data $(\mcN,\mcE,\mathcal C_\square)$ from \cite{continuumqft},
the coupling $g_3>0$, masses $m_F>0$ and $m_\phi>0$, and the distance
regularization $\eta>0$.  Then:
\begin{enumerate}[leftmargin=1.5em]
\item the promoted field manifold
\[
\mcX_\mathcal L:=SU(3)^{\mcE}\times \bbB^{\mcN}\times \Lambda_{\mathrm F}\times \R^{\mcN}
\]
admits a canonical product measure obtained from normalized Haar measure on
$SU(3)$, normalized Lebesgue measure on the color ball $\bbB\subset\C^3$,
finite Berezin measure in the fermionic directions, and Lebesgue measure in the
scalar directions;
\item the partition function
\[
Z=\int_{\mcX_\mathcal L}\exp\!\bigl(-S_{\mathrm W}(U)-S_{\mathrm F}(U,\bar\psi,\psi)-S_\phi(\phi)\bigr)\,\dd\mu
\]
exists as a finite complex number;
\item for any rooted spanning forest $\mcT$ of the underlying undirected graph,
every rooted gauge orbit meets the maximal-tree gauge slice exactly once, the
Faddeev--Popov determinant is identically $1$, and the gauge-fixed partition
function is an ordinary finite-dimensional integral over the cotree links;
\item integrating out the fermions yields the exact fermion determinant
$\det M_{\mathrm F}(U)$, while the scalar integral is exactly Gaussian with
covariance $(L_\eta+m_\phi^2 I)^{-1}$.
\end{enumerate}
\end{theorem}

\begin{remark}[Scope of the paper]
\label{rem:scope}
Theorem~\ref{thm:main} is a complete quantization theorem for a \emph{finite
lattice}.  The proofs below use only the finite graph
\[
(\mcN,\mcE,\mathcal C_\square)
\]
and the finite-dimensional gauge and matter actions fixed in the present
construction.  The
measure-theoretic statements therefore concern exact finite-dimensional
integrals, while the Lie-algebraic discussion in later sections concerns only
local coordinate descriptions of Haar measure near the identity.  In the
broader series, these finite-$N$ quantizations are then passed to the
stationary thermodynamic limit.  More precisely, \cite{convergence} proves the
weak convergence
\[
\pi_{N,1}\Rightarrow \pi_\infty,
\]
while \cite{continuumqft} formulates the corresponding regulator finite action
and continuum limit picture on $\pi_\infty$.
The same mean-field law carries the logarithmic Sobolev inequality proved in
\cite{convergence}, and Proposition~\ref{prop:thermo-lsi-transfer} below makes
explicit how that stiffness transfers to the finite product measures
$\pi_\infty^{\otimes k}$ supporting the thermodynamic finite-lattice
observables.  The continuum-indexed observable class used for that
thermodynamic presentation is fixed next.  In addition, the paper records the
distinguished localized Wilson sector whose sampled deterministic
representatives are exactly the localized Wilson functionals of
\cite{continuumqft}; this is the direct bridge from the exact finite regulator
theory constructed here to the continuum Yang--Mills functional identified in
that companion paper.
\end{remark}

\begin{definition}[Continuum-indexed local gauge-invariant plaquette algebra]
\label{def:continuum-plaquette-algebra}
Fix an orientation \(0\le \mu<\nu\le 3\), a compactly supported test function
\(\chi\in C_c(\R^4)\), and a bounded continuous class function
\(\Phi:SU(3)\to\C\), i.e.
\[
\Phi(VUV^{-1})=\Phi(U)
\qquad
\forall\,U,V\in SU(3).
\]
The corresponding continuum-indexed local gauge-invariant plaquette observable
is denoted by
\[
\mathcal O_{\chi,\mu\nu,\Phi}.
\]
Let \(\mathfrak A_{\mathrm{loc}}^{\mathrm{g.i.}}\) be the unital algebra
generated by finite linear combinations and finite products of such symbols.
\end{definition}

\begin{definition}[Canonical finite-regulator approximation map]
\label{def:continuum-plaquette-embedding}
For each regulator size \(N\), use the deterministic flat regulator data of
\cite{continuumqft}: the lattice spacing \(a_N\), the plaquette family
\(\mathcal C_\square\), and the plaquette base-point map
\[
P\mapsto x_P\in\R^4.
\]
For each orientation \(0\le\mu<\nu\le3\), let
\(\mathcal C_{\square,\mu\nu}\subset\mathcal C_\square\) be the coordinate
plaquettes of type \((\mu,\nu)\).  Define
\[
\iota_N\bigl(\mathcal O_{\chi,\mu\nu,\Phi}\bigr)(U)
:=
\sum_{P\in\mathcal C_{\square,\mu\nu}}
a_N^4\,\chi(x_P)\,\Phi\bigl(U(P)\bigr).
\]
Extend \(\iota_N\) linearly and multiplicatively to all of
\(\mathfrak A_{\mathrm{loc}}^{\mathrm{g.i.}}\).
\end{definition}

\begin{proposition}[Locality and boundedness of the canonical gauge-invariant observables]
\label{prop:continuum-plaquette-package}
The observable package of
Definitions~\ref{def:continuum-plaquette-algebra}
and~\ref{def:continuum-plaquette-embedding} has the following properties.
\begin{enumerate}[leftmargin=1.5em]
\item For every \(O\in\mathfrak A_{\mathrm{loc}}^{\mathrm{g.i.}}\),
\(\iota_N(O)\) is gauge invariant on the finite regulator lattice.
\item If \(O\) is generated by test functions \(\chi_1,\dots,\chi_m\), then
\(\iota_N(O)\) depends only on plaquettes whose base points lie in the compact
set \(\bigcup_{j=1}^m\supp\chi_j\), and hence only on the link variables
meeting that compact window.
\item For each generator \(\mathcal O_{\chi,\mu\nu,\Phi}\), there exists
\(C_{\chi,\Phi}<\infty\), independent of \(N\), such that
\[
\sup_{U\in\mcU}
\left|
\iota_N\bigl(\mathcal O_{\chi,\mu\nu,\Phi}\bigr)(U)
\right|
\le
C_{\chi,\Phi}.
\]
\item For every \(O\in\mathfrak A_{\mathrm{loc}}^{\mathrm{g.i.}}\), the
unnormalized Euclidean functional
\[
\int_{\mcX_\mathcal L}\iota_N(O)\,e^{-S}\,\dd\mu
\]
is finite.
\end{enumerate}
\end{proposition}

\begin{proof}
For a closed plaquette \(P\), a local gauge transformation
\(\Omega\) sends \(U(P)\) to \(\Omega(x_P)U(P)\Omega(x_P)^{-1}\), so gauge
invariance of each generator follows from the class-function property of
\(\Phi\).  Finite products and linear combinations preserve gauge invariance.

The locality statement is immediate from the definition of \(\iota_N\): only
plaquettes with base point \(x_P\in\supp\chi_j\) can contribute to the generator
associated with \(\chi_j\), and finite products involve only the union of the
corresponding compact supports.

Because \(\Phi\) is bounded and \(\chi\) is compactly supported, the
deterministic regulator geometry from \cite{continuumqft} implies that the
number of contributing \((\mu,\nu)\)-plaquettes is \(O(a_N^{-4})\) on every
fixed compact window, while each term carries the weight \(a_N^4\).  Hence the
sum is uniformly bounded in \(N\) and \(U\).

Finally, \(\iota_N(O)\) is bounded on the bosonic configuration space, so
finiteness of the displayed Euclidean functional is exactly the
unnormalized-integrability statement of
Theorem~\ref{thm:partition-finite}.
\end{proof}

\begin{corollary}[Nontrivial local gauge-invariant probe]
\label{cor:continuum-plaquette-nontrivial}
The algebra \(\mathfrak A_{\mathrm{loc}}^{\mathrm{g.i.}}\) is nontrivial.
\end{corollary}

\begin{proof}
Take any nonzero \(\chi\in C_c(\R^4)\), any \(0\le\mu<\nu\le3\), and
\[
\Phi(U):=1-\frac13\Ree\,\tr U.
\]
This \(\Phi\) is a bounded continuous class function on \(SU(3)\), and it is
not constant.  Therefore, for every regulator level \(N\) for which
\(\supp\chi\) contains at least one \((\mu,\nu)\)-plaquette base point,
the representative
\(\iota_N(\mathcal O_{\chi,\mu\nu,\Phi})\) is a nonconstant gauge-invariant
observable on the finite regulator.
\end{proof}

\begin{definition}[Distinguished localized Wilson family]
\label{def:wilson-density-generators}
For \(\chi\in C_c(\R^4)\), let \(\mathcal W_\chi\) denote the corresponding
continuum-indexed localized Wilson symbol.  Its canonical finite-regulator
representative is the gauge-invariant plaquette sum
\[
\jmath_N(\mathcal W_\chi)(U)
:=
\frac{6}{g_3^2}
\sum_{0\le\mu<\nu\le3}\sum_{P\in\mathcal C_{\square,\mu\nu}}
\chi(x_P)\left(1-\frac13\Ree\,\tr U(P)\right).
\]
\end{definition}

\begin{theorem}[Canonical Wilson generators recover the sampled deterministic local continuum limit]
\label{thm:wilson-generator-limit}
Let \(\chi\in C_c(\R^4)\).  For every regulator size \(N\),
\[
\jmath_N(\mathcal W_\chi)(U)
=
\frac{6}{g_3^2}
\sum_{0\le\mu<\nu\le3}\sum_{P\in\mathcal C_{\square,\mu\nu}}
\chi(x_P)\left(1-\frac13\Ree\,\tr U(P)\right).
\]
In particular, when the gauge field \(U\) is specialized to the deterministic
sampled plaquette holonomies \(U_3^{(N)}(P)\) of \cite{continuumqft}, one has
\[
\jmath_N(\mathcal W_\chi)(U_3^{(N)})
=
S_{\mathrm W,\chi}^{(N)},
\]
the localized Wilson functional of the companion continuum-gauge paper.
Consequently, the local continuum-limit theorem of \cite{continuumqft} gives
\[
\jmath_N(\mathcal W_\chi)(U_3^{(N)})
\longrightarrow
S_{\mathrm{YM},\chi}^{\mathrm{cont}}[F^\infty]
\]
along every admissible deterministic flat regulator family of that paper.
\end{theorem}

\begin{proof}
The displayed formula is exactly Definition~\ref{def:wilson-density-generators}.
Specializing to the sampled deterministic
plaquette holonomies of \cite{continuumqft} identifies the right-hand side with
the localized Wilson functional \(S_{\mathrm W,\chi}^{(N)}\).  The convergence
statement is therefore precisely the local continuum-limit theorem proved in
that paper.
\end{proof}

\begin{corollary}[Regulator-independence of the distinguished Wilson sector in the sampled deterministic presentation]
\label{cor:wilson-generator-independence}
Within the admissible deterministic flat regulator class of \cite{continuumqft},
the limit of \(\jmath_N(\mathcal W_\chi)(U_3^{(N)})\) depends only on
\(\chi\) and the continuum curvature profile \(F^\infty\), not on the chosen
admissible regulator exhaustion.
\end{corollary}

\begin{proof}
Theorem~\ref{thm:wilson-generator-limit} identifies every admissible-sequence
limit with the same explicit continuum quantity
\[
S_{\mathrm{YM},\chi}^{\mathrm{cont}}[F^\infty].
\]
Hence the limit is independent of the particular admissible deterministic flat
regulator family.
\end{proof}

\begin{corollary}[Sampled deterministic continuum Wilson functional on the distinguished local sector]
\label{cor:continuum-wilson-functional}
The assignment
\[
\omega_{\mathrm W}^\infty(\mathcal W_\chi)
:=
S_{\mathrm{YM},\chi}^{\mathrm{cont}}[F^\infty]
\]
defines a linear gauge-invariant continuum functional on the linear span of the
symbols \(\{\mathcal W_\chi:\chi\in C_c(\R^4)\}\) associated with the sampled
deterministic holonomy family of \cite{continuumqft}, and the
deterministic finite-regulator representatives
\(\jmath_N(\mathcal W_\chi)(U_3^{(N)})\) converge to that functional along every
admissible deterministic flat regulator family.
\end{corollary}

\begin{proof}
Linearity in \(\chi\) is immediate from the definition of
\(S_{\mathrm{YM},\chi}^{\mathrm{cont}}[F^\infty]\).  The convergence statement
is exactly Theorem~\ref{thm:wilson-generator-limit}.
\end{proof}

\section{Configuration space and measure}
\label{sec:config}

\subsection{Imported lattice data}

We summarize only the regulator-gauge ingredients from \cite{continuumqft} and
the explicit matter conventions adopted in this paper that are required for
quantization.

\begin{definition}[Imported lattice package]
\label{def:imported-package}
For each regulator size \(N\), the deterministic regulator construction
supplies the finite regulator graph and Wilson sector
\[
(\mcN,\mcE,\mathcal C_\square)=
(\mcN^{(N)},\mcE^{(N)},\mathcal C_\square^{(N)})
\]
and we equip that graph with the following finite ingredients:
\begin{enumerate}[leftmargin=1.5em]
\item finite site and oriented edge sets $\mcN,\mcE$;
\item for every plaquette $P\in\mathcal C_\square$, an oriented boundary
\[
\partial P=(e_1,e_2,e_3,e_4),\qquad e_k\in\mcE,
\]
with the ordered plaquette holonomy
\[
U(P):=U(e_1)U(e_2)U(e_3)U(e_4);
\]
\item the Wilson action
\[
S_{\mathrm W}(U):=\beta\sum_{P\in\mathcal C_\square}
\left(1-\frac13\Ree\,\tr U(P)\right),
\qquad
\beta:=\frac{6}{g_3^2};
\]
\item the matter actions
\[
S_{\mathrm F}(U,\bar\psi,\psi)
:=
\sum_{e\in\mcE}\bar\psi_{t(e)}
\left(\psi_{t(e)}-U(e)\psi_{s(e)}\right)
+m_F\sum_{x\in\mcN}\bar\psi_x\psi_x,
\]
\[
S_\phi(\phi)
:=
\frac12\sum_{e=(x\to y)\in\mcE}\frac{(\phi_y-\phi_x)^2}{\ell_e+\eta}
+\frac{m_\phi^2}{2}\sum_{x\in\mcN}\phi_x^2,
\]
where $\ell_e$ is the fixed edge length induced by the regulator geometry and
$\eta>0$ is the distance regularization.
\end{enumerate}
\end{definition}

\begin{remark}
The gauge sector above is exactly the regulator Wilson sector of
\cite{continuumqft}.  The only change made there is interpretive: the sampled
background links are promoted to independent coordinates on $SU(3)^{\mcE}$.
The fermionic and scalar terms are then fixed explicitly on that same finite
graph by the formulas displayed in Definition~\ref{def:imported-package}.
\end{remark}

\subsection{Field manifold}

\begin{definition}[Gauge-link manifold]
\label{def:gauge-manifold}
Let
\[
\mcU:=SU(3)^{\mcE}
=\{U:\mcE\to SU(3)\}.
\]
Because $\mcE$ is finite and $SU(3)$ is a compact smooth Lie group of real
dimension $8$, $\mcU$ is a compact smooth manifold of real dimension
$8|\mcE|$.
\end{definition}

\begin{definition}[Auxiliary color manifold]
\label{def:color-manifold}
Let
\[
\bbB:=\{z\in\C^3:\|z\|<1\},
\qquad
\mcCaux:=\bbB^{\mcN}.
\]
This is the product of the open unit balls suggested by the color normalization
bound $\|\mathbf c_x\|<1$ proved in \cite{continuumqft}.  Since no independent
color action was supplied in the companion paper, this sector will be assigned
a normalized product measure and decouples from the dynamical integrals below.
\end{definition}

\begin{definition}[Matter field domains]
\label{def:matter-domains}
Let
\[
\mcS:=\R^{\mcN}
\]
be the scalar field space.  For the fermionic sector, write
\[
V_{\mathrm F}:=\bigoplus_{x\in\mcN}\C_x^3
\]
and let $\Lambda_{\mathrm F}$ be the Grassmann algebra generated by
\[
\{\psi_{x,\alpha},\bar\psi_{x,\alpha}:x\in\mcN,\ \alpha=1,2,3\}.
\]
The full promoted field manifold is
\[
\mcX_\mathcal L:=\mcU\times\mcCaux\times \Lambda_{\mathrm F}\times\mcS.
\]
\end{definition}

\begin{remark}
Only the gauge and scalar sectors are manifolds in the ordinary sense.  The
fermionic sector is finite-dimensional Grassmann algebra, so integration in
those variables is algebraic rather than measure theoretic.  This is standard
for finite lattice fermions \cite{montvay1994quantum,rothe2012lattice}.
\end{remark}

\subsection{Measures}

\begin{definition}[Normalized Haar measure]
\label{def:haar}
Let $\mu_H$ denote the unique normalized Haar probability measure on $SU(3)$,
so that
\[
\mu_H(SU(3))=1
\]
and for every Borel set $A\subseteq SU(3)$ and every $g\in SU(3)$,
\[
\mu_H(gA)=\mu_H(Ag)=\mu_H(A).
\]
The gauge-link measure on $\mcU$ is the product probability measure
\[
\dd\mu_{\mcU}(U):=\prod_{e\in\mcE}\dd\mu_H(U(e)).
\]
\end{definition}

\begin{definition}[Auxiliary color measure]
\label{def:color-measure}
Let $\nu_{\bbB}$ be normalized Lebesgue measure on $\bbB\subset\C^3\cong\R^6$.
The color measure on $\mcCaux$ is
\[
\dd\nu_{\mcCaux}(c):=\prod_{x\in\mcN}\dd\nu_{\bbB}(c_x).
\]
By normalization, $\nu_{\mcCaux}(\mcCaux)=1$.
\end{definition}

\begin{definition}[Berezin and scalar measures]
\label{def:berezin}
The Berezin measure is
\[
\dd\mu_{\mathrm F}:=
\prod_{x\in\mcN}\prod_{\alpha=1}^3
\dd\bar\psi_{x,\alpha}\,\dd\psi_{x,\alpha},
\]
with the convention
\[
\int \dd\bar\psi\,\dd\psi\ e^{-\bar\psi\psi}=1.
\]
The scalar measure on $\mcS$ is ordinary Lebesgue measure
\[
\dd\phi:=\prod_{x\in\mcN}\dd\phi_x.
\]
\end{definition}

\begin{definition}[Total finite-lattice measure]
\label{def:total-measure}
The total measure symbol on $\mcX_\mathcal L$ is
\[
\dd\mu:=
\dd\mu_{\mcU}(U)\,\dd\nu_{\mcCaux}(c)\,\dd\mu_{\mathrm F}\,\dd\phi.
\]
When the integrand does not depend on $c$, the auxiliary color sector
contributes the unit factor
\[
\int_{\mcCaux}\dd\nu_{\mcCaux}(c)=1.
\]
\end{definition}

\begin{proposition}[Well-posedness of the finite measure factors]
\label{prop:measure-wellposed}
Each factor in Definition~\ref{def:total-measure} is well defined.  In
particular, $\mu_{\mcU}$ is a Borel probability measure on a compact manifold,
$\nu_{\mcCaux}$ is a Borel probability measure on a bounded open subset of
$\R^{6|\mcN|}$, $\mu_{\mathrm F}$ is a finite Berezin integration rule on a
finite Grassmann algebra, and $\phi\mapsto S_\phi(\phi)$ is a measurable
function on $\R^{|\mcN|}$.
\end{proposition}
\begin{proof}
The product group $SU(3)^{\mcE}$ is compact because $\mcE$ is finite and
$SU(3)$ is compact, hence the product Haar measure exists.  The set
$\mcCaux=\bbB^{\mcN}$ is a bounded open subset of a finite-dimensional Euclidean
space, so the normalized Lebesgue measure exists.  The fermionic sector has
finitely many generators, hence finite Berezin integration is an algebraic
operation.  Finally $S_\phi$ is a finite sum of continuous functions of
$\phi$, because each coefficient $(\ell_e+\eta)^{-1}$ is finite and positive.
\end{proof}

\section{Construction of the partition function}
\label{sec:partition}

\paragraph{The total Euclidean action.}
The quantized action on $\mcX_\mathcal L$ is the same algebraic expression as
in the companion lattice paper, but now interpreted on the promoted field
space:
\[
S(U,\bar\psi,\psi,\phi)
:=
S_{\mathrm W}(U)+S_{\mathrm F}(U,\bar\psi,\psi)+S_\phi(\phi).
\]
The auxiliary color variable $c\in\mcCaux$ does not enter this action and is
therefore normalized away.  No additional asymptotic limit is taken at this stage.
Every term is a finite sum over the fixed finite index sets $\mcN$, $\mcE$, and
$\mathcal C_\square$.

\paragraph{The master partition function.}
The central object is the finite-lattice partition function
\begin{equation}
\label{eq:partition-function}
Z
:=
\int_{\mcX_\mathcal L}
\exp\!\bigl(-S(U,\bar\psi,\psi,\phi)\bigr)\,\dd\mu.
\end{equation}
Because the color sector is normalized and decoupled,
\[
Z
=
\int_{\mcU\times\Lambda_{\mathrm F}\times\mcS}
\exp\!\bigl(-S_{\mathrm W}(U)-S_{\mathrm F}(U,\bar\psi,\psi)-S_\phi(\phi)\bigr)
\dd\mu_{\mcU}(U)\,\dd\mu_{\mathrm F}\,\dd\phi.
\]

\paragraph{Unnormalized Euclidean functionals.}
For any bounded gauge-invariant observable $\mcO$ whose fermionic degree is
even, define the unnormalized Euclidean functional by
\begin{equation}
\label{eq:unnormalized-expectation}
\widetilde\omega_N(\mcO)
:=
\int_{\mcX_\mathcal L}
\mcO(U,c,\bar\psi,\psi,\phi)\,
e^{-S(U,\bar\psi,\psi,\phi)}\,\dd\mu.
\end{equation}

\paragraph{Normalization and expectation values.}
When $Z\neq 0$, the normalized vacuum expectation is defined by
\begin{equation}
\label{eq:expectation}
\langle \mcO\rangle
:=
\frac{1}{Z}
\widetilde\omega_N(\mcO),
\end{equation}
provided $Z\neq 0$.  The nonvanishing of $Z$ is model dependent and cannot be
asserted for arbitrary finite complex fermion determinants.  What can always be
proved is finiteness of the numerator and denominator as complex numbers, which
is sufficient for all formal manipulations before division by $Z$.

\begin{theorem}[Existence and finiteness of the partition function]
\label{thm:partition-finite}
Under the hypotheses of Theorem~\ref{thm:main}, the partition function
\eqref{eq:partition-function} is a finite complex number.  Moreover, for every
bounded measurable observable $\mcO$ depending polynomially on the Grassmann
variables, the unnormalized expectation
\[
\int_{\mcX_\mathcal L}\mcO\,e^{-S}\,\dd\mu
\]
is finite.
\end{theorem}
\begin{proof}
Integrate the fermions first.  Section~\ref{sec:matter-integration} will show
that
\[
\int_{\Lambda_{\mathrm F}}e^{-S_{\mathrm F}(U,\bar\psi,\psi)}\,\dd\mu_{\mathrm F}
=\det M_{\mathrm F}(U),
\]
where $M_{\mathrm F}(U)$ is a finite complex matrix depending continuously on
$U$.  Since $\mcU$ is compact, the function $U\mapsto \det M_{\mathrm F}(U)$ is
bounded:
\[
\sup_{U\in\mcU}|\det M_{\mathrm F}(U)|<\infty.
\]
The Wilson term satisfies $S_{\mathrm W}(U)\ge 0$ for every $U$, because each
summand $1-\frac13\Ree\,\tr U(P)$ is nonnegative.  Hence
\[
|e^{-S_{\mathrm W}(U)}|\le 1.
\]
For the scalar sector,
\[
S_\phi(\phi)
\ge
\frac{m_\phi^2}{2}\sum_{x\in\mcN}\phi_x^2,
\]
because the edge term is nonnegative.
Therefore
\[
e^{-S_\phi(\phi)}
\le
\exp\!\left(-\frac{m_\phi^2}{2}\sum_{x\in\mcN}\phi_x^2\right),
\]
and the right-hand side is integrable on $\R^{|\mcN|}$.
Putting the bounds together,
\[
|Z|
\le
\sup_{U\in\mcU}|\det M_{\mathrm F}(U)|
\int_{\R^{|\mcN|}}
e^{-\frac{m_\phi^2}{2}\sum_x\phi_x^2}\,\dd\phi
<\infty.
\]
The same domination argument works after multiplying by any bounded observable
that is polynomial in the finitely many Grassmann generators, because Berezin
integration of such a polynomial again yields a bounded continuous function of
$U$.
\end{proof}

\section{Gauge redundancy and orbit analysis}
\label{sec:orbits}

\paragraph{The gauge orbit problem.}
The finite lattice theory has a local gauge symmetry at every site.  This means
that the coordinates $(U,\psi,\bar\psi,c)$ on field space are not physical
coordinates: many different tuples encode the same gauge state.  The redundancy
is continuous, because the gauge group is a product of copies of the compact Lie
group $SU(3)$.

\begin{definition}[Local gauge group]
\label{def:gauge-group}
The lattice gauge group is
\[
\mcG:=SU(3)^{\mcN}
=\{\Omega:\mcN\to SU(3)\}.
\]
Its action on the promoted fields is
\[
U^\Omega(e):=\Omega(t(e))\,U(e)\,\Omega(s(e))^{-1},
\]
\[
\psi_x^\Omega:=\Omega(x)\psi_x,\qquad
\bar\psi_x^\Omega:=\bar\psi_x\Omega(x)^{-1},\qquad
c_x^\Omega:=\Omega(x)c_x,\qquad
\phi_x^\Omega:=\phi_x.
\]
\end{definition}

\begin{proposition}[Gauge invariance of action and measure]
\label{prop:gauge-measure}
For every $\Omega\in\mcG$,
\[
S_{\mathrm W}(U^\Omega)=S_{\mathrm W}(U),\qquad
S_{\mathrm F}(U^\Omega,\bar\psi^\Omega,\psi^\Omega)=S_{\mathrm F}(U,\bar\psi,\psi),
\qquad
S_\phi(\phi^\Omega)=S_\phi(\phi).
\]
Moreover, $\mu_{\mcU}$, $\nu_{\mcCaux}$, and $\mu_{\mathrm F}$ are invariant
under the corresponding gauge transformations.  Consequently the total weighted
measure $e^{-S}\dd\mu$ is gauge invariant.
\end{proposition}
\begin{proof}
Gauge invariance of $S_{\mathrm W}$ was proved in \cite{continuumqft}.  For the
matter sector, a direct calculation from the defining formulas shows that
$S_{\mathrm F}$ is invariant under
\[
\psi_x\mapsto\Omega(x)\psi_x,\qquad
\bar\psi_x\mapsto\bar\psi_x\Omega(x)^{-1},
\]
because the gauge factors cancel pointwise in each edge term and in each mass
term, while $S_\phi$ is invariant because the scalar field is gauge neutral.
Haar invariance of $\mu_{\mcU}$
follows edgewise from left and right invariance of Haar measure on $SU(3)$.
The normalized Lebesgue measure on the ball $\bbB\subset\C^3$ is invariant
under the unitary action of $SU(3)$ because the action preserves Euclidean
norms and Jacobian $1$.  Berezin measure is invariant under the linear change of
variables $(\psi,\bar\psi)\mapsto(\Omega\psi,\bar\psi\Omega^{-1})$ because the
Jacobians $\det\Omega$ and $\det\Omega^{-1}$ cancel and $\det\Omega=1$ in
$SU(3)$.
\end{proof}

\begin{definition}[Gauge orbit and stabilizer]
\label{def:orbit}
For a field configuration
\[
\Phi:=(U,c,\bar\psi,\psi,\phi)\in\mcX_\mathcal L
\]
its gauge orbit is
\[
\mcO(\Phi):=\{\Phi^\Omega:\Omega\in\mcG\},
\]
and its stabilizer is
\[
\mathrm{Stab}(\Phi):=\{\Omega\in\mcG:\Phi^\Omega=\Phi\}.
\]
\end{definition}

\paragraph{Symmetry redundancy versus physical degrees of freedom.}
The physical configuration space is not $\mcX_\mathcal L$ itself but the orbit
space $\mcX_\mathcal L/\mcG$.  On a finite graph one can count the redundant and
nonredundant link degrees explicitly.  This count is most transparent after
choosing a spanning forest.

\begin{definition}[Underlying undirected graph and component count]
\label{def:components}
Let $|\mathcal L|$ be the undirected graph obtained from $(\mcN,\mcE)$ by
forgetting orientation and identifying parallel opposite orientations when both
are present.  Let $b_0(|\mathcal L|)$ be the number of connected components of
$|\mathcal L|$.
\end{definition}

\begin{proposition}[Dimension count for gauge redundancy]
\label{prop:dimension-count}
Assume the underlying undirected graph $|\mathcal L|$ has
$b_0:=b_0(|\mathcal L|)$ connected components.  Then:
\begin{enumerate}[leftmargin=1.5em]
\item the gauge-link manifold has real dimension $8|\mcE|$;
\item the rooted gauge group obtained by fixing one vertex per connected
component has real dimension $8(|\mcN|-b_0)$;
\item after maximal-tree gauge fixing, the number of remaining independent gauge
link coordinates is
\[
8\bigl(|\mcE|-|\mcN|+b_0\bigr),
\]
which is $8$ times the cycle rank of the graph.
\end{enumerate}
\end{proposition}
\begin{proof}
Item (1) is immediate from $\dim SU(3)=8$.  For item (2), each unfixed site
contributes one copy of $SU(3)$.  For item (3), every spanning forest on a
graph with $|\mcN|$ vertices and $b_0$ connected components has
$|\mcN|-b_0$ edges, so removing exactly those tree-edge coordinates leaves
$|\mcE|-|\mcN|+b_0$ cotree edges.  Each cotree edge carries $8$ real parameters.
\end{proof}

\begin{remark}
The compactness of $SU(3)$ means that the gauge-group probability measure is
finite when normalized.  Thus on a finite lattice the problem with gauge
redundancy is not an \emph{infinite numerical volume divergence}; the problem is
that the parametrization is many-to-one and therefore obscures the physical
degrees of freedom.  Gauge fixing removes this redundancy exactly.
\end{remark}

\section{Gauge fixing and the Faddeev--Popov procedure}
\label{sec:gauge-fixing}

\paragraph{Selecting the gauge condition.}
For a finite graph the cleanest globally valid gauge choice is a rooted
maximal-tree gauge.  Unlike lattice Landau gauge, it has no Gribov ambiguity:
every rooted orbit intersects the gauge slice exactly once.

\begin{definition}[Rooted spanning forest]
\label{def:tree}
Choose one root vertex in each connected component of $|\mathcal L|$ and let
$R\subseteq\mcN$ be the set of roots.  A \emph{rooted spanning forest}
$\mcT\subseteq\mcE$ is an oriented set of edges such that:
\begin{enumerate}[leftmargin=1.5em]
\item after forgetting orientation, $\mcT$ is a spanning forest of $|\mathcal L|$;
\item every edge of $\mcT$ is oriented away from its component root;
\item every nonroot vertex $x\in\mcN\setminus R$ has a unique parent edge
\[
e_x=(p(x)\to x)\in\mcT.
\]
\end{enumerate}
\end{definition}

\begin{definition}[Rooted gauge group and maximal-tree gauge slice]
\label{def:rooted-gauge}
The rooted gauge group is
\[
\mcG_*:=\{\Omega\in\mcG:\Omega(r)=I\ \text{for every root } r\in R\}.
\]
The maximal-tree gauge slice is
\[
\Sigma_{\mcT}:=\{U\in\mcU:U(e)=I\ \text{for every } e\in\mcT\}.
\]
\end{definition}

\begin{theorem}[Global rooted maximal-tree gauge]
\label{thm:tree-gauge}
For every $U\in\mcU$ there exists a unique $\Omega_U\in\mcG_*$ such that
\[
U^{\Omega_U}\in\Sigma_{\mcT}.
\]
Equivalently, every rooted gauge orbit in $\mcU$ intersects $\Sigma_{\mcT}$
exactly once.
\end{theorem}
\begin{proof}
Existence is recursive along each tree component.  Set $\Omega_U(r)=I$ at every
root.  Suppose $\Omega_U(p(x))$ is already defined for the parent of a nonroot
vertex $x$.  The condition
\[
U^{\Omega_U}(e_x)=I
\]
for the parent edge $e_x=(p(x)\to x)$ is
\[
\Omega_U(x)\,U(e_x)\,\Omega_U(p(x))^{-1}=I,
\]
hence
\[
\Omega_U(x)=\Omega_U(p(x))\,U(e_x)^{-1}.
\]
Because every nonroot vertex has a unique parent and every component contains a
unique path from the root to any vertex, this recursion defines $\Omega_U$
uniquely on all vertices.  By construction $\Omega_U\in\mcG_*$ and
$U^{\Omega_U}(e)=I$ for all $e\in\mcT$, so $U^{\Omega_U}\in\Sigma_{\mcT}$.

For uniqueness, let $\Omega_1,\Omega_2\in\mcG_*$ satisfy
$U^{\Omega_1},U^{\Omega_2}\in\Sigma_{\mcT}$.  Set
$\Gamma:=\Omega_2^{-1}\Omega_1\in\mcG_*$.  Then for every tree edge
$e=(x\to y)\in\mcT$,
\[
I
=
\bigl(U^{\Omega_1}\bigr)(e)
=
\Gamma(y)\,\bigl(U^{\Omega_2}\bigr)(e)\,\Gamma(x)^{-1}
=
\Gamma(y)\Gamma(x)^{-1}.
\]
Thus $\Gamma(y)=\Gamma(x)$ on every tree edge.  Since $\Gamma(r)=I$ on every
root and every vertex is connected to a root by a unique tree path, induction
along the tree gives $\Gamma(x)=I$ for all $x$.  Hence $\Omega_1=\Omega_2$.
\end{proof}

\paragraph{The Faddeev--Popov identity.}
Because the gauge slice is global and unique, the Faddeev--Popov formula is
particularly simple.

\begin{definition}[Delta distribution on the tree gauge slice]
\label{def:delta-tree}
Let $\delta_I$ be the Dirac distribution at the identity of $SU(3)$, defined by
\[
\int_{SU(3)}f(V)\,\delta_I(V)\,\dd\mu_H(V)=f(I)
\]
for continuous $f$.  Define the tree-gauge delta distribution by
\[
\delta_{\mcT}(U):=\prod_{e\in\mcT}\delta_I\!\bigl(U(e)\bigr).
\]
\end{definition}

\begin{theorem}[Exact Faddeev--Popov identity in maximal-tree gauge]
\label{thm:fp-identity}
For every $U\in\mcU$,
\[
1=\int_{\mcG_*}\delta_{\mcT}(U^\Omega)\,\dd\mu_{\mcG_*}(\Omega),
\]
where $\mu_{\mcG_*}$ is product normalized Haar measure on $\mcG_*$.  Hence for
every gauge-invariant integrable functional $F$ on $\mcU$,
\[
\int_{\mcU}F(U)\,\dd\mu_{\mcU}(U)
=
\int_{\mcU}F(U)\,\delta_{\mcT}(U)\,\dd\mu_{\mcU}(U).
\]
\end{theorem}
\begin{proof}
Fix $U\in\mcU$.  We prove the first identity by induction on the number of
nonroot vertices.
If the rooted forest has no nonroot vertices, then $\mcT=\varnothing$ and
$\delta_{\mcT}\equiv 1$, so the claim is immediate.
Assume now that $\mcT$ has at least one edge and choose a leaf
$x\in\mcN\setminus R$.
Because $x$ has no children, the only factor in $\delta_{\mcT}(U^\Omega)$ that
depends on $\Omega(x)$ is the factor associated with the parent edge
$e_x=(p(x)\to x)$:
\[
\delta_I\!\bigl(U^\Omega(e_x)\bigr)
=
\delta_I\!\bigl(\Omega(x)U(e_x)\Omega(p(x))^{-1}\bigr).
\]
For fixed values of all other gauge variables, set
$g_x:=U(e_x)\Omega(p(x))^{-1}\in SU(3)$.  Then left invariance of Haar measure
gives
\[
\int_{SU(3)}\delta_I\!\bigl(\Omega(x)g_x\bigr)\,\dd\mu_H(\Omega(x))
=
\int_{SU(3)}\delta_I(V)\,\dd\mu_H(V)=1.
\]
Thus integrating over the leaf variable contributes exactly the unit factor and
removes the leaf edge from the product.
Deleting the leaf and repeating the same argument inductively over the smaller
rooted forest proves
\[
\int_{\mcG_*}\delta_{\mcT}(U^\Omega)\,\dd\mu_{\mcG_*}(\Omega)=1.
\]

For the second identity, insert the first into the integral of $F$:
\[
\int_{\mcU}F(U)\,\dd\mu_{\mcU}(U)
=
\int_{\mcU}\int_{\mcG_*}
F(U)\,\delta_{\mcT}(U^\Omega)\,\dd\mu_{\mcG_*}(\Omega)\,\dd\mu_{\mcU}(U).
\]
Use Fubini's theorem and then change variables $U\mapsto U^{\Omega^{-1}}$.
Haar invariance implies $\dd\mu_{\mcU}(U^{\Omega^{-1}})=\dd\mu_{\mcU}(U)$, and
gauge invariance of $F$ gives $F(U^{\Omega^{-1}})=F(U)$.  Thus
\[
\int_{\mcU}F(U)\,\dd\mu_{\mcU}(U)
=
\int_{\mcU}F(U)\,\delta_{\mcT}(U)\,\dd\mu_{\mcU}(U)
\int_{\mcG_*}\dd\mu_{\mcG_*}(\Omega).
\]
The last factor equals $1$ because $\mu_{\mcG_*}$ is normalized.
\end{proof}

\paragraph{The Faddeev--Popov determinant.}
The determinant can be computed directly from infinitesimal gauge variations.
In maximal-tree gauge it is constant.

\begin{definition}[Infinitesimal gauge parameters]
\label{def:inf-gauge}
Let
\[
T^a:=\frac{\lambda_a}{2},\qquad a=1,\dots,8,
\]
where $\lambda_a$ are the Gell-Mann matrices.  For each site $x\in\mcN$ let
\[
\omega_x=\sum_{a=1}^8\omega_x^a T^a\in\mcHh,
\]
and define the infinitesimal rooted gauge transformation
\[
\Omega_\varepsilon(x):=\exp(i\varepsilon\,\omega_x),
\qquad \omega_r=0\ \text{for every }r\in R.
\]
\end{definition}

\begin{proposition}[Linearized gauge map on the tree]
\label{prop:linearized-fp}
Let $U\in\Sigma_{\mcT}$, so $U(e)=I$ for all $e\in\mcT$.  For a tree edge
$e=(x\to y)\in\mcT$,
\[
-i\left.\frac{\dd}{\dd\varepsilon}\right|_{\varepsilon=0}
U^{\Omega_\varepsilon}(e)
=
\omega_y-\omega_x.
\]
In the basis $\{T^a\}_{a=1}^8$, the linearized Faddeev--Popov operator is
\[
M_{\mathrm{FP}}^{(\mcT)} = B_{\mcT}\otimes I_8,
\]
where $B_{\mcT}$ is the rooted incidence matrix of the spanning forest.
\end{proposition}
\begin{proof}
Differentiate
\[
U^{\Omega_\varepsilon}(e)
=
\Omega_\varepsilon(y)\,U(e)\,\Omega_\varepsilon(x)^{-1}
\]
at $\varepsilon=0$.  Since $U(e)=I$ on the slice,
\[
-i\left.\frac{\dd}{\dd\varepsilon}\right|_{\varepsilon=0}
U^{\Omega_\varepsilon}(e)
=
\omega_y-\omega_x.
\]
Writing each $\omega_x$ in the Gell-Mann basis shows that the eight color
components decouple and each is acted on by the same incidence matrix
$B_{\mcT}$.
\end{proof}

\begin{proposition}[The Faddeev--Popov determinant is identically one]
\label{prop:fp-one}
With the rooted maximal-tree gauge choice,
\[
\Delta_{\mathrm{FP}}^{(\mcT)}=\det M_{\mathrm{FP}}^{(\mcT)}=1.
\]
\end{proposition}
\begin{proof}
Order the nonroot vertices so that every parent precedes every child.  Index
the rows of $B_{\mcT}$ by the tree edges $e_x=(p(x)\to x)$ and the columns by
the nonroot vertices $x$.  Then the row corresponding to $e_x$ has entry
$+1$ in column $x$, entry $-1$ in column $p(x)$ if $p(x)$ is nonroot, and zero
in all columns associated with vertices not on the path from the root to $x$.
By the chosen ordering this matrix is triangular with diagonal entries all equal
to $1$.  Hence
\[
\det B_{\mcT}=1.
\]
Therefore
\[
\det M_{\mathrm{FP}}^{(\mcT)}
=
\det(B_{\mcT}\otimes I_8)
=
(\det B_{\mcT})^8
=
1.
\]
\end{proof}

\begin{corollary}[Gauge-fixed partition function]
\label{cor:gauge-fixed-z}
The partition function may be written as
\[
Z=
\int_{\mcU\times\Lambda_{\mathrm F}\times\mcS}
\delta_{\mcT}(U)\,
e^{-S_{\mathrm W}(U)-S_{\mathrm F}(U,\bar\psi,\psi)-S_\phi(\phi)}
\dd\mu_{\mcU}(U)\,\dd\mu_{\mathrm F}\,\dd\phi.
\]
After integrating out the constrained tree links, this becomes an ordinary
integral over the cotree edge variables
\[
\mcE_{\mathrm{co}}:=\mcE\setminus\mcT.
\]
\end{corollary}
\begin{proof}
Apply Theorem~\ref{thm:fp-identity} to the gauge-invariant integrand
obtained after combining the matter variables with
Proposition~\ref{prop:gauge-measure}.  Since $\delta_{\mcT}(U)$ forces
$U(e)=I$ on all tree edges, integrating over those edges contributes the unit
factor
\[
\int_{SU(3)}\delta_I(V)\,\dd\mu_H(V)=1
\]
edge by edge.
\end{proof}

\section{Ghost fields and the effective action}
\label{sec:ghosts}

\paragraph{Introducing ghost fields.}
Although the Faddeev--Popov determinant is trivial in rooted maximal-tree
gauge, it is still useful to represent it as a finite Berezin integral, because
this makes the gauge-fixing algebra parallel to the standard continuum
formalism.  The continuum-indexed gauge-invariant observable package fixed in
Definitions~\ref{def:continuum-plaquette-algebra}
and~\ref{def:continuum-plaquette-embedding} is independent of these auxiliary
variables; the present section records the exact finite-dimensional algebraic
repackaging of the same gauge-fixed measure.

\begin{definition}[Ghost and antighost fields]
\label{def:ghosts}
For each nonroot vertex $x\in\mcN\setminus R$, let
\[
c_x=\sum_{a=1}^8 c_x^a T^a
\]
be an anticommuting $\mcHh$-valued ghost field.  For each tree edge
$e\in\mcT$, let
\[
\bar c_e=\sum_{a=1}^8 \bar c_e^a T^a
\]
be an anticommuting $\mcHh$-valued antighost field.  Their Berezin measure is
\[
\dd\mu_{\mathrm{gh}}
:=
\prod_{e\in\mcT}\prod_{a=1}^8 \dd\bar c_e^a
\prod_{x\in\mcN\setminus R}\prod_{a=1}^8 \dd c_x^a.
\]
For every root vertex $r\in R$, set
\[
c_r:=0.
\]
\end{definition}

\begin{definition}[Ghost action in maximal-tree gauge]
\label{def:ghost-action}
Define
\[
S_{\mathrm{ghost}}^{(\mcT)}(\bar c,c)
:=
2\sum_{e=(x\to y)\in\mcT}\tr\!\bigl(\bar c_e(c_y-c_x)\bigr),
\]
with the convention $c_r:=0$ for root vertices $r\in R$.
\end{definition}

\begin{proposition}[Ghost representation of the Faddeev--Popov determinant]
\label{prop:ghost-det}
The ghost integral reproduces the Faddeev--Popov determinant:
\[
\int e^{-S_{\mathrm{ghost}}^{(\mcT)}(\bar c,c)}\,\dd\mu_{\mathrm{gh}}
=\Delta_{\mathrm{FP}}^{(\mcT)}=1.
\]
\end{proposition}
\begin{proof}
Expand the ghost fields in the basis $\{T^a\}_{a=1}^8$.  Then the action
factorizes over color indices:
\[
S_{\mathrm{ghost}}^{(\mcT)}(\bar c,c)
=
\sum_{a=1}^8\sum_{e=(x\to y)\in\mcT}
\bar c_e^a(c_y^a-c_x^a).
\]
Here we used $\tr(T^aT^b)=\frac12\delta_{ab}$.
This is precisely the quadratic Berezin form associated with the matrix
$B_{\mcT}\otimes I_8$, so the standard finite-dimensional Berezin identity gives
\[
\int e^{-S_{\mathrm{ghost}}^{(\mcT)}}\,\dd\mu_{\mathrm{gh}}
=
\det(B_{\mcT}\otimes I_8).
\]
Proposition~\ref{prop:fp-one} gives the value $1$.
\end{proof}

\paragraph{The effective quantized action.}
To display the standard BRST structure it is convenient to introduce auxiliary
bosonic Nakanishi--Lautrup fields $b_e\in\mcHh$ for $e\in\mcT$.

\begin{definition}[BRST differential and gauge-fixing fermion]
\label{def:brst}
Let $s$ act on the full field algebra by
\[
s\,U(e)=i c_{t(e)}U(e)-iU(e)c_{s(e)},
\qquad
s\,\psi_x=i c_x\psi_x,
\qquad
s\,\bar\psi_x=-i\bar\psi_x c_x,
\]
\[
s\,\phi_x=0,
\qquad
s\,c_x=-i c_x^2,
\qquad
s\,\bar c_e=b_e,
\qquad
s\,b_e=0.
\]
Here the rooted ghost convention from Definition~\ref{def:ghosts} is in force,
so $c_r=0$ for every root $r\in R$.
Define the gauge-fixing fermion
\[
\Psi_{\mcT}
:=
\sum_{e\in\mcT}
2\tr\!\left(\bar c_e\Bigl(-i\bigl(U(e)-I\bigr)+\frac{\alpha}{2}b_e\Bigr)\right),
\qquad \alpha\ge 0.
\]
\end{definition}

\begin{lemma}[Nilpotency of the BRST differential]
\label{lem:brst}
The operator $s$ satisfies $s^2=0$ on generators and therefore on the whole
field algebra.
\end{lemma}
\begin{proof}
For $\phi_x$ and $b_e$ the claim is immediate.  For $\bar c_e$,
$s^2\bar c_e=s b_e=0$.  For $c_x$,
\[
s^2 c_x = -i\,s(c_x^2) = -i\bigl((s c_x)c_x-c_x(s c_x)\bigr)
=-i\bigl((-ic_x^2)c_x-c_x(-ic_x^2)\bigr)=0.
\]
For $\psi_x$,
\[
s^2\psi_x=i(s c_x)\psi_x-i c_x(s\psi_x)
=i(-ic_x^2)\psi_x-i c_x(i c_x\psi_x)=0,
\]
and similarly for $\bar\psi_x$.  Finally, for a link variable $U(e)$,
\begin{align*}
s^2 U(e)
&=
i(s c_{t(e)})U(e)+i c_{t(e)}sU(e)-i sU(e)c_{s(e)}-iU(e)s c_{s(e)}\\
&=
c_{t(e)}^2U(e)-c_{t(e)}^2U(e)+U(e)c_{s(e)}^2-U(e)c_{s(e)}^2=0.
\end{align*}
\end{proof}

\begin{proposition}[BRST-exact gauge fixing]
\label{prop:brst-exact}
The gauge-fixing and ghost terms are
\[
S_{\mathrm{gf+gh}}^{(\mcT)}:=s\Psi_{\mcT}.
\]
Explicitly,
\[
S_{\mathrm{gf+gh}}^{(\mcT)}
=
\sum_{e\in\mcT}2\tr\!\left(
-i\,b_e\bigl(U(e)-I\bigr)+\frac{\alpha}{2}b_e^2
+\bar c_e\bigl(c_{t(e)}U(e)-U(e)c_{s(e)}\bigr)
\right).
\]
Consequently the effective action
\[
S_{\mathrm{eff}}
:=
S_{\mathrm W}+S_{\mathrm F}+S_\phi+S_{\mathrm{gf+gh}}^{(\mcT)}
\]
is BRST invariant:
\[
sS_{\mathrm{eff}}=0.
\]
\end{proposition}
\begin{proof}
Apply $s$ to $\Psi_{\mcT}$ and use Definition~\ref{def:brst}.  Since the
fields $b_e$ commute with everything and the ghosts anticommute, one obtains
the displayed formula.  Gauge invariance of
$S_{\mathrm W}+S_{\mathrm F}+S_\phi$ implies $s(S_{\mathrm W}+S_{\mathrm F}+S_\phi)=0$,
because $s$ is the infinitesimal gauge differential with ghost parameter $c$.
Then
\[
sS_{\mathrm{eff}}
=
s(S_{\mathrm W}+S_{\mathrm F}+S_\phi)+s^2\Psi_{\mcT}=0
\]
by Lemma~\ref{lem:brst}.
\end{proof}

\begin{remark}
On the maximal-tree gauge slice $U(e)=I$ for $e\in\mcT$, the ghost term in
$S_{\mathrm{gf+gh}}^{(\mcT)}$ reduces exactly to
$S_{\mathrm{ghost}}^{(\mcT)}$ from Definition~\ref{def:ghost-action}.  Since
the Faddeev--Popov determinant is constant, the ghost sector decouples after
gauge fixing.  The BRST presentation is therefore an algebraic repackaging of
the same finite gauge-fixing data.
\end{remark}

\section{Matter sector integration}
\label{sec:matter-integration}

\paragraph{Fermionic Berezin integration.}
On a finite lattice the fermion integral is exact and elementary once the
quadratic form is written in matrix notation.

\begin{definition}[Fermion matrix]
\label{def:fermion-matrix}
Order the site-color basis vectors by pairs $(x,\alpha)$ with
$x\in\mcN$ and $\alpha\in\{1,2,3\}$.  Define the matrix
$M_{\mathrm F}(U)\in\C^{3|\mcN|\times 3|\mcN|}$ by
\[
\bigl(M_{\mathrm F}(U)\bigr)_{(y,\alpha),(x,\beta)}
:=
\left(m_F+\indeg(y)\right)\delta_{xy}\delta_{\alpha\beta}
-\sum_{\substack{e\in\mcE\\ s(e)=x,\ t(e)=y}} U(e)_{\alpha\beta},
\]
where
\[
\indeg(y):=\#\{e\in\mcE:t(e)=y\}.
\]
\end{definition}

\begin{lemma}[Quadratic form identity]
\label{lem:fermion-quadratic}
With $M_{\mathrm F}(U)$ from Definition~\ref{def:fermion-matrix},
\[
S_{\mathrm F}(U,\bar\psi,\psi)
=
\sum_{x,y\in\mcN}\sum_{\alpha,\beta=1}^3
\bar\psi_{y,\alpha}\,
\bigl(M_{\mathrm F}(U)\bigr)_{(y,\alpha),(x,\beta)}\,
\psi_{x,\beta}.
\]
\end{lemma}
\begin{proof}
Expand the edge sum in Definition~\ref{def:imported-package}:
\[
\sum_{e\in\mcE}\bar\psi_{t(e)}\psi_{t(e)}
=
\sum_{y\in\mcN}\indeg(y)\,\bar\psi_y\psi_y.
\]
The remaining transported terms contribute the off-diagonal blocks
$-\sum_{e:s(e)=x,t(e)=y}U(e)$.  Adding the mass term gives the stated matrix.
\end{proof}

\begin{theorem}[Fermion determinant]
\label{thm:fermion-det}
For every gauge field $U\in\mcU$,
\[
\int_{\Lambda_{\mathrm F}}
e^{-S_{\mathrm F}(U,\bar\psi,\psi)}\,\dd\mu_{\mathrm F}
=
\det M_{\mathrm F}(U).
\]
Consequently
\[
Z
=
\int_{\mcU\times\mcS}
e^{-S_{\mathrm W}(U)-S_\phi(\phi)}\det M_{\mathrm F}(U)\,
\dd\mu_{\mcU}(U)\,\dd\phi.
\]
\end{theorem}
\begin{proof}
By Lemma~\ref{lem:fermion-quadratic}, the exponent is the standard quadratic
Berezin form $\bar\Psi M_{\mathrm F}(U)\Psi$ in finitely many Grassmann
variables.  The finite-dimensional Berezin integral identity therefore gives
\[
\int e^{-\bar\Psi M\Psi}\,\dd\bar\Psi\,\dd\Psi=\det M
\]
for any finite matrix $M$.  Applying this with $M=M_{\mathrm F}(U)$ yields the
first claim.  Substituting the result into the definition of $Z$ gives the
second claim.
\end{proof}

\begin{remark}[Virtual loops and graph support]
\label{rem:fermion-loops}
The determinant $\det M_{\mathrm F}(U)$ is a finite polynomial in the link
matrices and encodes the usual closed virtual-fermion loops.  Because the link
set $\mcE$ is the finite edge set from the deterministic regulator package
\cite{continuumqft}, only loops on this explicit graph contribute.  The
path support is therefore inherited entirely from that regulator incidence
structure, with no additional construction-specific combinatorics appearing in this
finite-dimensional determinant expansion.
\end{remark}

\paragraph{Scalar kinetic fluctuations.}
The scalar action is purely bosonic, real valued, and free massive.  The
regularization $\eta>0$ excludes edge-length singularities in the kinetic
coefficients.

\begin{definition}[Scalar stiffness operator]
\label{def:scalar-operator}
Define the symmetric quadratic form
\[
Q_\eta(\phi)
:=
\frac12\sum_{e=(x\to y)\in\mcE}\frac{(\phi_y-\phi_x)^2}{\ell_e+\eta}.
\]
There is a unique symmetric matrix $L_\eta\in\R^{|\mcN|\times|\mcN|}$ such that
\[
Q_\eta(\phi)=\frac12\,\phi^\top L_\eta \phi.
\]
\end{definition}

\begin{lemma}[Positivity and bounded coefficients]
\label{lem:scalar-positive}
The operator $L_\eta$ is positive semidefinite.  Moreover every edge weight
\[
w_e:=\frac{1}{\ell_e+\eta}
\]
satisfies
\[
0<w_e\le \eta^{-1}.
\]
\end{lemma}
\begin{proof}
Each summand in $Q_\eta(\phi)$ is nonnegative, so $Q_\eta(\phi)\ge 0$ for every
$\phi$, hence $L_\eta$ is positive semidefinite.  Since $\ell_e\ge 0$ and
$\eta>0$,
\[
\ell_e+\eta\ge \eta,
\]
which implies the upper bound on $w_e$.
\end{proof}

\begin{theorem}[Scalar integrability]
\label{thm:scalar-integrable}
\[
Z_\phi
:=
\int_{\mcS}e^{-S_\phi(\phi)}\,\dd\phi
\]
is finite, and
\[
Z_\phi
=
(2\pi)^{|\mcN|/2}
\det(L_\eta+m_\phi^2 I)^{-1/2}.
\]
\end{theorem}
\begin{proof}
By Definition~\ref{def:scalar-operator},
\[
S_\phi(\phi)=\frac12\phi^\top L_\eta\phi
+\frac{m_\phi^2}{2}\phi^\top\phi.
\]
Since $L_\eta$ is positive semidefinite and $m_\phi^2 I$ is strictly positive,
the matrix $L_\eta+m_\phi^2 I$ is positive definite.  The standard Gaussian
integral formula gives the displayed determinant expression.
\end{proof}

\begin{proposition}[Scalar source integral]
\label{prop:scalar-source}
For every source vector $J=(J_x)_{x\in\mcN}\in\R^{|\mcN|}$,
\[
Z_\phi[J]
:=
\int_{\mcS}\exp\!\left(-S_\phi(\phi)+\sum_{x\in\mcN}J_x\phi_x\right)\dd\phi
\]
is finite and satisfies
\[
Z_\phi[J]
=
Z_\phi[0]\,
\exp\!\left(\frac12\,J^\top (L_\eta+m_\phi^2 I)^{-1}J\right).
\]
In particular, all scalar polynomial moments are finite.
\end{proposition}
\begin{proof}
Write
\[
K_\phi:=L_\eta+m_\phi^2 I.
\]
By Theorem~\ref{thm:scalar-integrable}, the matrix $K_\phi$ is positive
definite.  Completing the square gives
\[
\frac12\phi^\top K_\phi\phi-J^\top\phi
=
\frac12(\phi-K_\phi^{-1}J)^\top K_\phi(\phi-K_\phi^{-1}J)
-\frac12 J^\top K_\phi^{-1}J.
\]
Substituting this identity into the integral and translating the integration
variable by $K_\phi^{-1}J$ yields the stated formula.  Differentiating the
right-hand side with respect to the source components proves finiteness of all
scalar polynomial moments.
\end{proof}

\begin{remark}[Role of $\eta$]
\label{rem:eta}
The parameter $\eta>0$ is not a perturbative convenience; it is the exact
mechanism that prevents the scalar kinetic coefficients from diverging when two
sites become spacetime-coincident.  Without $\eta$, the coefficient
$1/\ell_e$ would be singular on zero-length edges.  With $\eta>0$, every
kinetic coefficient is uniformly bounded by $\eta^{-1}$.
\end{remark}

\section{Lie-algebra coordinates and local measure structure}
\label{sec:local-measure}

\paragraph{Local exponential coordinates.}
The finite-lattice path integral is defined globally by Haar measure on each
copy of $SU(3)$.  To compare that measure with the usual Lie-algebra variables
used in perturbation theory and in local small-field expansions, one works in
an exponential chart around the identity.

\begin{definition}[Symmetric exponential chart]
\label{def:exp-chart}
Let $\mcHh$ be the Hermitian traceless $3\times 3$ matrices with inner product
\[
\langle H,K\rangle:=\frac12\tr(HK).
\]
Choose an open neighborhood $\mathcal V\subset\mcHh$ of $0$ such that:
\begin{enumerate}[leftmargin=1.5em]
\item $\mathcal V$ is symmetric: $H\in\mathcal V$ implies $-H\in\mathcal V$;
\item the map
\[
\exp_i:\mathcal V\to SU(3),
\qquad
\exp_i(H):=e^{iH},
\]
is a diffeomorphism onto an open neighborhood $\mathcal U$ of the identity.
\end{enumerate}
\end{definition}

\begin{proposition}[Single-edge Haar density in exponential coordinates]
\label{prop:haar-density}
Let $\mu_H$ be normalized Haar measure on $SU(3)$ and let
$\mathcal V,\mathcal U$ be as in Definition~\ref{def:exp-chart}.  Then there
exists a smooth strictly positive function
\[
J:\mathcal V\to(0,\infty)
\]
such that for every continuous function $f$ supported in $\mathcal U$,
\[
\int_{SU(3)} f(U)\,\dd\mu_H(U)
=
\int_{\mathcal V} f\!\bigl(e^{iH}\bigr)\,J(H)\,\dd H,
\]
where $\dd H$ is Lebesgue measure on the $8$-dimensional Euclidean space
$\mcHh$.  After fixing the normalization of $\dd H$ so that $J(0)=1$, the
density is even and satisfies
\[
J(H)=1+O(\|H\|^2)
\qquad\text{as }H\to 0.
\]
\end{proposition}
\begin{proof}
Push forward $\mu_H|_{\mathcal U}$ by the inverse chart
$\log_i:=\exp_i^{-1}:\mathcal U\to\mathcal V$.
Because $\exp_i$ is a diffeomorphism, the pushforward is a smooth measure on
the Euclidean open set $\mathcal V$, so it has a smooth strictly positive
density $J(H)$ with respect to Lebesgue measure.
Since Lebesgue measure on $\mcHh$ is unique only up to a positive constant,
we may normalize it so that $J(0)=1$.

The inversion map on $SU(3)$ sends $e^{iH}$ to $e^{-iH}$, and Haar measure is
invariant under inversion.  Because $\mathcal V$ is symmetric, this implies
\[
J(H)=J(-H)
\qquad\text{for all }H\in\mathcal V.
\]
Hence the first-order term in the Taylor expansion of $J$ at the origin
vanishes, and smoothness gives
\[
J(H)=J(0)+O(\|H\|^2)=1+O(\|H\|^2).
\]
\end{proof}

\begin{corollary}[Second-order matching with Lebesgue measure]
\label{cor:haar-second-order}
In the chart of Definition~\ref{def:exp-chart},
\[
\dd\mu_H\!\bigl(e^{iH}\bigr)=\bigl(1+O(\|H\|^2)\bigr)\dd H
\qquad\text{as }H\to 0.
\]
Equivalently, for every bounded set $B\subset\mcHh$ and every $a>0$ small
enough,
\[
\dd\mu_H\!\bigl(e^{i aA}\bigr)
=
a^8\bigl(1+O(a^2)\bigr)\dd A
\]
uniformly for $A\in B$.
\end{corollary}
\begin{proof}
The first statement is a restatement of Proposition~\ref{prop:haar-density}.
For the second, substitute $H=aA$.  Since $\mcHh$ has real dimension $8$,
\[
\dd H=a^8\,\dd A,
\]
and for $A$ in a fixed bounded set the estimate
\[
J(aA)=1+O(a^2\|A\|^2)=1+O(a^2)
\]
is uniform.
\end{proof}

\begin{remark}[Thermodynamic role of the local chart]
The local coordinate description of Haar measure is a fixed-$N$ statement on a
single edge.  In the broader thermodynamic programme it supplies the local
measure model for the finite-lattice link and plaquette observables that appear in
the $N\to\infty$ limit described in \cite{convergence,continuumqft}.
\end{remark}

\begin{proposition}[Transfer of the mean-field LSI to thermodynamic observables]
\label{prop:thermo-lsi-transfer}
Let $\pi_\infty$ be the stationary mean-field law from \cite{convergence}.
Assume, as proved there, that there exist constants $a,b,c\in\R$ with
$ac>b^2$ and a constant $C_{\mathrm{LSI}}<\infty$ such that
\[
\Ent_{\pi_\infty}(g^2)
\le
C_{\mathrm{LSI}}
\int_{\R^8}\Gamma_{\mathrm{hypo}}^\infty(g,g)\,\dd\pi_\infty
\]
for every smooth compactly supported $g:\R^8\to\R$, where
\[
\Gamma_{\mathrm{hypo}}^\infty(g,g)
:=
a\|\nabla_v g\|^2
+2b\,\nabla_x g\cdot\nabla_v g
+c\|\nabla_x g\|^2.
\]
For each $k\ge1$, let
\[
\Pi_k^\infty:=\pi_\infty^{\otimes k},
\]
and define
\[
\Gamma_{\mathrm{hypo}}^{(k)}(F,F)
:=
\sum_{j=1}^k \Gamma_{\mathrm{hypo},j}^\infty(F,F),
\]
where $\Gamma_{\mathrm{hypo},j}^\infty$ acts on the $j$-th variable
$z_j=(x_j,v_j)\in\R^8$.  Then
\[
\Ent_{\Pi_k^\infty}(F^2)
\le
C_{\mathrm{LSI}}
\int_{(\R^8)^k}\Gamma_{\mathrm{hypo}}^{(k)}(F,F)\,\dd\Pi_k^\infty
\]
for every smooth compactly supported $F:(\R^8)^k\to\R$.
In particular, the finite-lattice link observables and plaquette observables of
\cite{continuumqft}, which are evaluated on $\pi_\infty^{\otimes2}$ and
$\pi_\infty^{\otimes4}$, are supported on equilibrium product measures
carrying the same logarithmic Sobolev constant $C_{\mathrm{LSI}}$.
\end{proposition}
\begin{proof}
We argue by induction on $k$.  The case $k=1$ is exactly the mean-field
logarithmic Sobolev inequality proved in \cite{convergence}.

Assume the claim has been proved for $k-1$, and write
\[
z=(z',z_k)\in (\R^8)^{k-1}\times\R^8.
\]
For a smooth compactly supported $F:(\R^8)^k\to\R$, set
\[
F_{z_k}(z'):=F(z',z_k),
\qquad
m(z_k):=\int_{(\R^8)^{k-1}} F(z',z_k)^2\,\dd\Pi_{k-1}^\infty(z'),
\qquad
G(z_k):=m(z_k)^{1/2}.
\]
The entropy chain rule for product measures gives
\[
\Ent_{\Pi_k^\infty}(F^2)
=
\int_{\R^8}\Ent_{\Pi_{k-1}^\infty}(F_{z_k}^2)\,\dd\pi_\infty(z_k)
+\Ent_{\pi_\infty}(G^2).
\]
Applying the induction hypothesis to the first term yields
\[
\int_{\R^8}\Ent_{\Pi_{k-1}^\infty}(F_{z_k}^2)\,\dd\pi_\infty(z_k)
\le
C_{\mathrm{LSI}}
\sum_{j=1}^{k-1}
\int_{(\R^8)^k}\Gamma_{\mathrm{hypo},j}^\infty(F,F)\,\dd\Pi_k^\infty.
\]

For the second term, regularize the square root by
\[
G_\varepsilon(z_k):=(m(z_k)+\varepsilon)^{1/2},
\qquad \varepsilon>0.
\]
Because $F$ is smooth and compactly supported, the function $m$ is smooth and
bounded, hence $G_\varepsilon$ is smooth and bounded on $\R^8$.  Let
\[
M:=
\begin{pmatrix}
cI_4 & bI_4\\
bI_4 & aI_4
\end{pmatrix},
\]
so that
\[
\Gamma_{\mathrm{hypo},k}^\infty(F,F)
=
\left\langle \nabla_{z_k}F,\,
M\,\nabla_{z_k}F\right\rangle,
\qquad
\nabla_{z_k}:=(\nabla_{x_k},\nabla_{v_k}).
\]
Because $ac>b^2$, the matrix $M$ is positive definite.  On the set where
$m(z_k)>0$,
\[
\nabla_{z_k}G_\varepsilon(z_k)
=
\frac{1}{G_\varepsilon(z_k)}
\int_{(\R^8)^{k-1}}F(z',z_k)\,\nabla_{z_k}F(z',z_k)\,\dd\Pi_{k-1}^\infty(z').
\]
Therefore
\begin{align*}
G_\varepsilon(z_k)^2\,\Gamma_{\mathrm{hypo}}^\infty(G_\varepsilon,G_\varepsilon)
&=
\left\langle
\int F\,\nabla_{z_k}F\,\dd\Pi_{k-1}^\infty,\,
M
\int F\,\nabla_{z_k}F\,\dd\Pi_{k-1}^\infty
\right\rangle\\
&\le
\left(\int F^2\,\dd\Pi_{k-1}^\infty\right)
\left(
\int
\left\langle \nabla_{z_k}F,\,
M\,\nabla_{z_k}F\right\rangle
\dd\Pi_{k-1}^\infty
\right)\\
&=
m(z_k)
\int_{(\R^8)^{k-1}}
\Gamma_{\mathrm{hypo},k}^\infty(F,F)(z',z_k)\,
\dd\Pi_{k-1}^\infty(z').
\end{align*}
Hence
\[
\Gamma_{\mathrm{hypo}}^\infty(G_\varepsilon,G_\varepsilon)(z_k)
\le
\frac{m(z_k)}{m(z_k)+\varepsilon}
\int_{(\R^8)^{k-1}}
\Gamma_{\mathrm{hypo},k}^\infty(F,F)(z',z_k)\,
\dd\Pi_{k-1}^\infty(z')
\le
\]
\[
\int_{(\R^8)^{k-1}}
\Gamma_{\mathrm{hypo},k}^\infty(F,F)(z',z_k)\,
\dd\Pi_{k-1}^\infty(z').
\]
To invoke the one-particle inequality from \cite{convergence}, choose
$\chi_R\in C_c^\infty(\R^8)$ such that
\[
0\le \chi_R\le 1,\qquad
\chi_R\equiv 1\text{ on }B_R,\qquad
\operatorname{supp}\chi_R\subset B_{2R},\qquad
\|\nabla \chi_R\|_\infty\le \frac{C_\chi}{R}.
\]
Set
\[
H_{\varepsilon,R}:=\chi_R G_\varepsilon\in C_c^\infty(\R^8).
\]
Applying the one-particle inequality to $H_{\varepsilon,R}$ gives
\[
\Ent_{\pi_\infty}(H_{\varepsilon,R}^2)
\le
C_{\mathrm{LSI}}
\int_{\R^8}\Gamma_{\mathrm{hypo}}^\infty(H_{\varepsilon,R},H_{\varepsilon,R})\,
\dd\pi_\infty.
\]
Since $G_\varepsilon$ is bounded and $\chi_R\to 1$ pointwise, dominated
convergence gives
\[
\Ent_{\pi_\infty}(H_{\varepsilon,R}^2)\to \Ent_{\pi_\infty}(G_\varepsilon^2)
\qquad\text{as }R\to\infty.
\]
Moreover,
\[
\nabla H_{\varepsilon,R}
=
\chi_R\nabla G_\varepsilon+G_\varepsilon\nabla\chi_R,
\]
so positivity of $M$ implies
\[
\Gamma_{\mathrm{hypo}}^\infty(H_{\varepsilon,R},H_{\varepsilon,R})
\le
2\chi_R^2\Gamma_{\mathrm{hypo}}^\infty(G_\varepsilon,G_\varepsilon)
+
2G_\varepsilon^2\Gamma_{\mathrm{hypo}}^\infty(\chi_R,\chi_R).
\]
Because $a,b,c$ are fixed, there exists a constant $C_M<\infty$ such that
\[
\Gamma_{\mathrm{hypo}}^\infty(\chi_R,\chi_R)
\le
\frac{C_M}{R^2}\one_{B_{2R}\setminus B_R}.
\]
Hence the second term converges to $0$ in $L^1(\pi_\infty)$, while the first
term converges by dominated convergence, using the pointwise bound above by the
$\Pi_{k-1}^\infty$-average of $\Gamma_{\mathrm{hypo},k}^\infty(F,F)$.  Passing
to the limit $R\to\infty$ yields
\[
\Ent_{\pi_\infty}(G_\varepsilon^2)
\le
C_{\mathrm{LSI}}
\int_{\R^8}\Gamma_{\mathrm{hypo}}^\infty(G_\varepsilon,G_\varepsilon)\,\dd\pi_\infty.
\]
Integrating the previous pointwise estimate then gives
\[
\Ent_{\pi_\infty}(G_\varepsilon^2)
\le
C_{\mathrm{LSI}}
\int_{(\R^8)^k}\Gamma_{\mathrm{hypo},k}^\infty(F,F)\,\dd\Pi_k^\infty.
\]
Since $G_\varepsilon^2=m+\varepsilon$ decreases pointwise to $G^2=m$ as
$\varepsilon\downarrow0$, and $m$ is bounded while $t\mapsto t\log t$ is
continuous on $[0,\infty)$, dominated convergence yields
\[
\Ent_{\pi_\infty}(G_\varepsilon^2)\to \Ent_{\pi_\infty}(G^2).
\]
Combining the bounds for the two entropy terms proves the claim for $k$.
\end{proof}

\begin{remark}[Indicator-weighted thermodynamic observables]
The thermodynamic link and plaquette observables of \cite{continuumqft} are
integrated against bounded indicators $\chi_t$ and $\chi_\square$ on the
product spaces $(\R^8)^2$ and $(\R^8)^4$.  Since
\[
0\le \chi_t,\chi_\square\le 1,
\]
their integrability and entropy control are inherited from the supporting
product measures $\pi_\infty^{\otimes2}$ and $\pi_\infty^{\otimes4}$ of
Proposition~\ref{prop:thermo-lsi-transfer}; no additional hypothesis is
required at this stage.
\end{remark}

\begin{remark}[Compatibility with thermodynamic product measures]
\label{rem:thermo-observable-class}
Let \(O\in\mathfrak A_{\mathrm{loc}}^{\mathrm{g.i.}}\) be a product of \(m\)
generators from Definition~\ref{def:continuum-plaquette-algebra}.  In the
thermodynamic plaquette presentation of \cite{continuumqft}, an \(m\)-fold
product of single-plaquette gauge-invariant observables is naturally evaluated
on a \(4m\)-point equilibrium product space.  Proposition~\ref{prop:thermo-lsi-transfer}
therefore applies to the supporting product measure \(\pi_\infty^{\otimes 4m}\),
with the same logarithmic Sobolev inequality constant \(C_{\mathrm{LSI}}\) as
for \(\pi_\infty\).
\end{remark}

\section{Discussion of excitations and propagators}
\label{sec:propagators}

\paragraph{The propagator on the finite regulator graph.}
Because the theory is finite, all correlation functions can be defined by
source differentiation exactly.  Their graph-theoretic support is determined by
the finite regulator edge graph.

\begin{definition}[Generating functional]
\label{def:generating}
Introduce scalar sources $J=(J_x)_{x\in\mcN}$ and fermionic sources
$\eta=(\eta_x)_{x\in\mcN}$, $\bar\eta=(\bar\eta_x)_{x\in\mcN}$.  The generating
functional is
\[
Z[J,\eta,\bar\eta]
:=
\int
\exp\!\left(
-S
+\sum_{x\in\mcN}J_x\phi_x
+\sum_{x\in\mcN}\bar\eta_x\psi_x
+\sum_{x\in\mcN}\bar\psi_x\eta_x
\right)\dd\mu.
\]
\end{definition}

\begin{definition}[Two-point functions]
\label{def:twopoint}
Whenever $Z[0,0,0]\neq 0$, the scalar and fermion two-point functions are
\[
G_\phi(x,y):=\langle \phi_x\phi_y\rangle,
\qquad
S_F(x,y)_{\alpha\beta}:=\langle \psi_{x,\alpha}\bar\psi_{y,\beta}\rangle.
\]
\end{definition}

\begin{proposition}[Background-field propagators]
\label{prop:background-prop}
Fix a gauge field $U$.
\begin{enumerate}[leftmargin=1.5em]
\item The scalar background covariance is
\[
G_\phi^{(U)}=(L_\eta+m_\phi^2 I)^{-1}.
\]
\item The fermion background
covariance is
\[
S_F^{(U)}=M_{\mathrm F}(U)^{-1}.
\]
\end{enumerate}
\end{proposition}
\begin{proof}
The scalar claim is the standard Gaussian covariance formula from
Theorem~\ref{thm:scalar-integrable}.  By
Proposition~\ref{prop:fermion-invertible}, the matrix $M_{\mathrm F}(U)$ is
invertible for every $U$.  The fermion claim then follows from differentiating
the Berezin generating functional
\[
\int e^{-\bar\Psi M_{\mathrm F}(U)\Psi+\bar\eta\Psi+\bar\Psi\eta}\,\dd\bar\Psi\,\dd\Psi
=
\det M_{\mathrm F}(U)\,
e^{\bar\eta M_{\mathrm F}(U)^{-1}\eta}.
\]
\end{proof}

\begin{corollary}[Exact finite Euclidean functionals on the canonical local gauge-invariant sectors]
\label{cor:finite-euclidean-local-sectors}
For every regulator size \(N\), every
\(O\in\mathfrak A_{\mathrm{loc}}^{\mathrm{g.i.}}\), and every
\(\chi\in C_c(\R^4)\), the unnormalized Euclidean insertions
\[
\widetilde\omega_N\bigl(\iota_N(O)\bigr)
\qquad\text{and}\qquad
\widetilde\omega_N\bigl(\jmath_N(\mathcal W_\chi)\bigr)
\]
are finite complex numbers.
\end{corollary}

\begin{proof}
By Proposition~\ref{prop:continuum-plaquette-package}, each
\(\iota_N(O)\) is a bounded gauge-invariant bosonic observable.  By
Definition~\ref{def:wilson-density-generators}, each
\(\jmath_N(\mathcal W_\chi)\) is a finite linear combination of bounded
plaquette class functions, hence is again a bounded gauge-invariant bosonic
observable.  The finiteness of both unnormalized Euclidean functionals is
therefore an immediate application of Theorem~\ref{thm:partition-finite}.
\end{proof}

\begin{theorem}[Exact finite bosonic reduction of gauge-invariant insertions]
\label{thm:bosonic-reduction}
Fix a rooted spanning forest \(\mcT\) as in
Theorem~\ref{thm:tree-gauge}.  Let \(\mcO:\mcU\to\C\) be a bounded measurable
gauge-invariant observable, viewed as an observable on
\(\mcX_{\mathcal L}\) by ignoring the color, fermion, and scalar variables.
Then the unnormalized Euclidean functional satisfies
\[
\widetilde\omega_N(\mcO)
=
Z_\phi
\int_{\mcU}
\mcO(U)\,e^{-S_{\mathrm W}(U)}\det M_{\mathrm F}(U)\,
\dd\mu_{\mcU}(U),
\]
where
\[
Z_\phi=(2\pi)^{|\mcN|/2}\det(L_\eta+m_\phi^2 I)^{-1/2}.
\]
Moreover, equivalently,
\[
\widetilde\omega_N(\mcO)
=
Z_\phi
\int_{\mcU}
\delta_{\mcT}(U)\,\mcO(U)\,e^{-S_{\mathrm W}(U)}\det M_{\mathrm F}(U)\,
\dd\mu_{\mcU}(U),
\]
and after integrating out the constrained tree links this becomes an ordinary
finite-dimensional integral over the cotree edge variables
\(\mcE_{\mathrm{co}}=\mcE\setminus\mcT\).
\end{theorem}

\begin{proof}
Because the color sector is normalized and decoupled from the action, the
color integration contributes the unit factor.  Since \(\mcO\) depends only on
the gauge field, \eqref{eq:unnormalized-expectation} gives
\[
\widetilde\omega_N(\mcO)
=
\int_{\mcU\times\Lambda_{\mathrm F}\times\mcS}
\mcO(U)\,e^{-S_{\mathrm W}(U)-S_{\mathrm F}(U,\bar\psi,\psi)-S_\phi(\phi)}
\dd\mu_{\mcU}(U)\,\dd\mu_{\mathrm F}\,\dd\phi.
\]
Apply Theorem~\ref{thm:fermion-det} to integrate out the fermions:
\[
\widetilde\omega_N(\mcO)
=
\int_{\mcU\times\mcS}
\mcO(U)\,e^{-S_{\mathrm W}(U)-S_\phi(\phi)}\det M_{\mathrm F}(U)\,
\dd\mu_{\mcU}(U)\,\dd\phi.
\]
Now apply Theorem~\ref{thm:scalar-integrable}.  Since the scalar action is
independent of \(U\), the scalar integral factors out as \(Z_\phi\), yielding
\[
\widetilde\omega_N(\mcO)
=
Z_\phi
\int_{\mcU}
\mcO(U)\,e^{-S_{\mathrm W}(U)}\det M_{\mathrm F}(U)\,
\dd\mu_{\mcU}(U).
\]

By Proposition~\ref{prop:gauge-measure}, the original finite action and product
measure are gauge invariant.  Since \(\mcO\) is gauge invariant, the full
integrand before the fermion and scalar integrations is gauge invariant.  Hence
the gauge-fixed representation from Corollary~\ref{cor:gauge-fixed-z} applies,
and because the Faddeev--Popov determinant is identically \(1\) by
Theorem~\ref{thm:fp-identity}, one obtains
\[
\widetilde\omega_N(\mcO)
=
Z_\phi
\int_{\mcU}
\delta_{\mcT}(U)\,\mcO(U)\,e^{-S_{\mathrm W}(U)}\det M_{\mathrm F}(U)\,
\dd\mu_{\mcU}(U).
\]
The final statement is exactly the last sentence of
Corollary~\ref{cor:gauge-fixed-z}.
\end{proof}

\begin{corollary}[Exact finite bosonic reduction on the canonical local gauge-invariant sectors]
\label{cor:bosonic-reduction-local-sectors}
For every regulator size \(N\), every
\(O\in\mathfrak A_{\mathrm{loc}}^{\mathrm{g.i.}}\), and every
\(\chi\in C_c(\R^4)\),
\begin{align*}
\widetilde\omega_N\bigl(\iota_N(O)\bigr)
&=
Z_\phi
\int_{\mcU}
\delta_{\mcT}(U)\,\iota_N(O)(U)\,e^{-S_{\mathrm W}(U)}\det M_{\mathrm F}(U)\,
\dd\mu_{\mcU}(U),\\
\widetilde\omega_N\bigl(\jmath_N(\mathcal W_\chi)\bigr)
&=
Z_\phi
\int_{\mcU}
\delta_{\mcT}(U)\,\jmath_N(\mathcal W_\chi)(U)\,e^{-S_{\mathrm W}(U)}
\det M_{\mathrm F}(U)\,\dd\mu_{\mcU}(U).
\end{align*}
In each case, after integrating out the constrained tree links, the right-hand
side is an ordinary finite-dimensional integral over the cotree edge variables.
\end{corollary}

\begin{proof}
By Proposition~\ref{prop:continuum-plaquette-package},
\(\iota_N(O)\) is gauge invariant for every
\(O\in\mathfrak A_{\mathrm{loc}}^{\mathrm{g.i.}}\).  By
Definition~\ref{def:wilson-density-generators},
\(\jmath_N(\mathcal W_\chi)\) is a finite linear combination of gauge-invariant
plaquette class functions and is therefore gauge invariant.  Apply
Theorem~\ref{thm:bosonic-reduction} to each observable.
\end{proof}

\begin{proposition}[Unconditional fermion invertibility]
\label{prop:fermion-invertible}
For every gauge field $U\in\mcU$, the matrix $M_{\mathrm F}(U)$ is invertible.
More precisely, if
\[
M_{\mathrm F}(U)=D-H(U),
\qquad
D:=\diag\bigl(m_F+\indeg(x)\bigr)_{x\in\mcN}\otimes I_3,
\]
then
\[
\|D^{-1}H(U)\|_{\infty,2\to\infty,2}
\le
\max_{y\in\mcN}\frac{\indeg(y)}{m_F+\indeg(y)}
<1,
\]
where
\[
\|\xi\|_{\infty,2}:=\max_{x\in\mcN}\|\xi_x\|_2
\qquad
\text{for }
\xi=(\xi_x)_{x\in\mcN}\in\bigoplus_{x\in\mcN}\C_x^3.
\]
\end{proposition}
\begin{proof}
For $\xi=(\xi_x)_{x\in\mcN}$, the $y$-block of $H(U)\xi$ is
\[
\bigl(H(U)\xi\bigr)_y
=
\sum_{\substack{e\in\mcE\\ t(e)=y}}U(e)\,\xi_{s(e)}.
\]
Since each $U(e)\in SU(3)$ is unitary,
\[
\|U(e)\xi_{s(e)}\|_2=\|\xi_{s(e)}\|_2.
\]
Therefore
\[
\bigl\|\bigl(H(U)\xi\bigr)_y\bigr\|_2
\le
\sum_{\substack{e\in\mcE\\ t(e)=y}}\|\xi_{s(e)}\|_2
\le
\indeg(y)\,\|\xi\|_{\infty,2}.
\]
Multiplying by $D^{-1}$ gives
\[
\bigl\|\bigl(D^{-1}H(U)\xi\bigr)_y\bigr\|_2
\le
\frac{\indeg(y)}{m_F+\indeg(y)}\,\|\xi\|_{\infty,2}.
\]
Taking the maximum over $y$ yields the stated bound.
Because $m_F>0$, every fraction $\indeg(y)/(m_F+\indeg(y))$ is strictly less
than $1$.  Hence $\|D^{-1}H(U)\|_{\infty,2\to\infty,2}<1$, so the Neumann
series for $(I-D^{-1}H(U))^{-1}$ converges.
Since
\[
M_{\mathrm F}(U)=D\bigl(I-D^{-1}H(U)\bigr)
\]
and $D$ is invertible, $M_{\mathrm F}(U)$ is invertible.
\end{proof}

\begin{proposition}[Fermion hopping expansion]
\label{prop:hopping}
Let
\[
M_{\mathrm F}(U)=D-H(U),
\]
where
\[
D:=\diag\bigl(m_F+\indeg(x)\bigr)_{x\in\mcN}\otimes I_3
\]
and $H(U)$ is the off-diagonal hopping part from
Proposition~\ref{prop:fermion-invertible}.  Then
\[
M_{\mathrm F}(U)^{-1}
=
\sum_{n=0}^\infty (D^{-1}H(U))^n D^{-1},
\]
and the $(x,y)$ matrix element is a convergent sum over directed paths
$\gamma:y\rightsquigarrow x$ in $\mcE$.  Therefore the support of the fermion
propagator is controlled by the finite regulator graph, and its
combinatorics are inherited from the deterministic regulator incidence of the
companion data in \cite{continuumqft}.
\end{proposition}
\begin{proof}
By Proposition~\ref{prop:fermion-invertible},
\[
\|D^{-1}H(U)\|_{\infty,2\to\infty,2}<1.
\]
Hence the Neumann series for $(I-D^{-1}H(U))^{-1}$ converges in that operator
norm, giving
\[
\bigl(I-D^{-1}H(U)\bigr)^{-1}
=
\sum_{n=0}^\infty (D^{-1}H(U))^n.
\]
Multiplying on the right by $D^{-1}$ yields the displayed formula for
$M_{\mathrm F}(U)^{-1}$.
Each power $(D^{-1}H(U))^n$ is a sum over length-$n$ compositions
of off-diagonal hopping operators, hence over directed paths of length $n$ in
the graph.  The only admissible steps for those paths are edges in $\mcE$, as
determined by the finite regulator data from \cite{continuumqft}.
\end{proof}

\begin{theorem}[Absolutely convergent fermion loop expansion and exact effective gauge action]
\label{thm:fermion-loop-expansion}
For every gauge field \(U\in\mcU\), define
\[
K(U):=D^{-1}H(U),
\qquad
q_*:=\max_{y\in\mcN}\frac{\indeg(y)}{m_F+\indeg(y)}<1.
\]
Then the following hold.
\begin{enumerate}[leftmargin=1.5em]
\item The spectrum of \(K(U)\) lies in the closed disk \(\{z\in\C:|z|\le q_*\}\).
\item The principal matrix logarithm of \(I-K(U)\) is well defined and satisfies
\[
\Log\!\bigl(I-K(U)\bigr)
=
-\sum_{n=1}^\infty \frac{K(U)^n}{n},
\]
with absolute convergence in the operator norm
\(\|\cdot\|_{\infty,2\to\infty,2}\).
\item The fermion determinant admits the exact logarithmic expansion
\[
\log\det M_{\mathrm F}(U)
=
\log\det D
-\sum_{n=1}^\infty \frac1n\,\Tr\!\bigl(K(U)^n\bigr).
\]
\item For each \(n\ge1\), let \(\mathcal L_n^{\mathrm{cl}}\) be the finite set
of closed based directed walks
\[
\ell=(e_1,\dots,e_n)
\]
in \(\mcE\), meaning
\[
t(e_j)=s(e_{j+1})\quad(1\le j\le n-1),
\qquad
t(e_n)=s(e_1).
\]
For \(\ell\in\mathcal L_n^{\mathrm{cl}}\), define
\[
w(\ell):=\prod_{j=1}^n \frac{1}{m_F+\indeg(t(e_j))},
\qquad
U(\ell):=U(e_n)\cdots U(e_1).
\]
Then
\[
\Tr\!\bigl(K(U)^n\bigr)
=
\sum_{\ell\in\mathcal L_n^{\mathrm{cl}}}
w(\ell)\,\tr\!\bigl(U(\ell)\bigr).
\]
In particular, each \(\Tr(K(U)^n)\) is gauge invariant.
\item The series is uniformly absolutely summable in the sense that
\[
\left|\Tr\!\bigl(K(U)^n\bigr)\right|
\le
3|\mcN|\,q_*^n,
\]
and therefore, for every \(L\ge0\),
\[
\sum_{n=L+1}^\infty \frac1n
\left|\Tr\!\bigl(K(U)^n\bigr)\right|
\le
\frac{3|\mcN|}{1-q_*}\,q_*^{L+1}.
\]
\end{enumerate}
\end{theorem}

\begin{proof}
By Proposition~\ref{prop:fermion-invertible},
\[
\|K(U)\|_{\infty,2\to\infty,2}
\le q_*<1.
\]
Hence the spectral radius satisfies
\[
\rho(K(U))\le \|K(U)\|_{\infty,2\to\infty,2}\le q_*,
\]
which proves the first claim.

Because \(\|K(U)\|_{\infty,2\to\infty,2}<1\), the scalar power series
\[
-\sum_{n=1}^\infty \frac{z^n}{n}
\]
for \(\log(1-z)\) converges absolutely on the closed disk
\(|z|\le q_*\).  Therefore the matrix series
\[
L(U):=-\sum_{n=1}^\infty \frac{K(U)^n}{n}
\]
converges absolutely in \(\|\cdot\|_{\infty,2\to\infty,2}\).  Since
\(\sigma(K(U))\subset\{|z|\le q_*\}\), the spectrum of \(I-K(U)\) is contained
in the disk \(\{w:|w-1|\le q_*\}\), which avoids the nonpositive real axis
because \(q_*<1\).  Thus the principal matrix logarithm is defined, and by the
holomorphic functional calculus for the principal branch,
\[
\Log\!\bigl(I-K(U)\bigr)=L(U).
\]
This proves the second claim.

Since
\[
M_{\mathrm F}(U)=D\bigl(I-K(U)\bigr),
\]
we have
\[
\det M_{\mathrm F}(U)=\det D\cdot \det\!\bigl(I-K(U)\bigr).
\]
The matrix \(D\) is diagonal with strictly positive real entries, so
\(\log\det D\) is the ordinary real logarithm.  In finite dimension,
\[
\det\!\bigl(I-K(U)\bigr)
=
\exp\!\Bigl(\Tr \Log\!\bigl(I-K(U)\bigr)\Bigr).
\]
Taking logarithms and inserting the series from the second claim yields
\[
\log\det M_{\mathrm F}(U)
=
\log\det D+\Tr \Log\!\bigl(I-K(U)\bigr)
=
\log\det D-\sum_{n=1}^\infty \frac1n\,\Tr\!\bigl(K(U)^n\bigr),
\]
which proves the third claim.

To identify the trace terms, note that for \(\xi=(\xi_x)_{x\in\mcN}\),
\[
\bigl(K(U)\xi\bigr)_y
=
\frac{1}{m_F+\indeg(y)}
\sum_{\substack{e\in\mcE\\ t(e)=y}}U(e)\,\xi_{s(e)}.
\]
Iterating this identity shows that the \((x,x)\)-block of \(K(U)^n\) is the
sum over all directed length-\(n\) paths
\[
x=x_0\xrightarrow{e_1}x_1\xrightarrow{e_2}\cdots\xrightarrow{e_n}x_n=x
\]
of the matrices
\[
\left(\prod_{j=1}^n \frac{1}{m_F+\indeg(x_j)}\right)
U(e_n)\cdots U(e_1).
\]
Taking the \(3\times3\) trace of the diagonal block and then summing over
\(x\in\mcN\) gives
\[
\Tr\!\bigl(K(U)^n\bigr)
=
\sum_{\ell\in\mathcal L_n^{\mathrm{cl}}}
w(\ell)\,\tr\!\bigl(U(\ell)\bigr),
\]
where \(\mathcal L_n^{\mathrm{cl}}\) is precisely the set of closed based
directed walks of length \(n\).  This proves the displayed loop expansion.

Under a gauge transformation \(U(e)\mapsto \Omega(t(e))U(e)\Omega(s(e))^{-1}\),
the product along a closed walk transforms by conjugation:
\[
U(\ell)\mapsto \Omega(s(e_1))\,U(\ell)\,\Omega(s(e_1))^{-1}.
\]
Hence \(\tr(U(\ell))\) is gauge invariant, and therefore so is
\(\Tr(K(U)^n)\).

Finally, for any matrix \(A\) on \(\bigoplus_{x\in\mcN}\C_x^3\),
\[
|\Tr A|
\le
\sum_{x\in\mcN} |\tr A_{xx}|
\le
3|\mcN|\,\|A\|_{\infty,2\to\infty,2}.
\]
Applying this with \(A=K(U)^n\) gives
\[
\left|\Tr\!\bigl(K(U)^n\bigr)\right|
\le
3|\mcN|\,\|K(U)\|_{\infty,2\to\infty,2}^n
\le
3|\mcN|\,q_*^n.
\]
Therefore
\[
\sum_{n=L+1}^\infty \frac1n
\left|\Tr\!\bigl(K(U)^n\bigr)\right|
\le
3|\mcN|\sum_{n=L+1}^\infty q_*^n
=
\frac{3|\mcN|}{1-q_*}\,q_*^{L+1},
\]
which proves the last claim.
\end{proof}

\begin{corollary}[Exact loop-effective bosonic action]
\label{cor:loop-effective-bosonic-action}
Define the exact fermionic effective gauge action by
\[
S_{\mathrm{eff},F}(U)
:=
\sum_{n=1}^\infty \frac1n\,\Tr\!\bigl(K(U)^n\bigr),
\qquad
K(U)=D^{-1}H(U).
\]
Then the series converges absolutely for every \(U\in\mcU\), each partial sum
is gauge invariant, and
\[
\det M_{\mathrm F}(U)
=
\det D\,e^{-S_{\mathrm{eff},F}(U)}.
\]
Consequently, for every bounded measurable gauge-invariant observable
\(\mcO:\mcU\to\C\),
\[
\widetilde\omega_N(\mcO)
=
Z_\phi\,\det D
\int_{\mcU}
\delta_{\mcT}(U)\,\mcO(U)\,
e^{-S_{\mathrm W}(U)-S_{\mathrm{eff},F}(U)}\,
\dd\mu_{\mcU}(U).
\]
\end{corollary}

\begin{proof}
Absolute convergence and gauge invariance follow directly from
Theorem~\ref{thm:fermion-loop-expansion}.  The determinant identity is the
exponential form of the logarithmic expansion proved there:
\[
\log\det M_{\mathrm F}(U)=\log\det D-S_{\mathrm{eff},F}(U).
\]
Exponentiating gives
\[
\det M_{\mathrm F}(U)=\det D\,e^{-S_{\mathrm{eff},F}(U)}.
\]
Substituting this identity into Theorem~\ref{thm:bosonic-reduction} yields the
stated bosonic reduction formula.
\end{proof}

\paragraph{Wilson loops and the area law.}
The Wilson loop remains the basic gauge-invariant probe of confinement.

\begin{definition}[Wilson loop expectation]
\label{def:wilson-loop-exp}
For an oriented closed loop $\gamma$ in the lattice,
\[
W(\gamma):=\tr\!\left(\prod_{e\subset \gamma}U(e)\right),
\qquad
\langle W(\gamma)\rangle := Z^{-1}\int W(\gamma)e^{-S}\,\dd\mu,
\]
whenever $Z\neq 0$.
\end{definition}

\begin{remark}[Area-law criterion]
\label{rem:area-law}
If a family of loops $\gamma$ admits spanning surfaces $\Sigma_\gamma$ in the
embedded spacetime geometry and one can prove
\[
|\langle W(\gamma)\rangle|
\le
C\,e^{-\sigma\,\mathrm{Area}(\Sigma_\gamma)}
\]
with $\sigma>0$ independent of the loop size, then one has the usual Euclidean
area-law signature of confinement \cite{wilson1974confinement}.  The present
paper does \emph{not} establish such a bound; it merely provides the exact
finite-lattice measure in which the question is well posed.
\end{remark}

\section{Conclusion}
\label{sec:conclusion}

\paragraph{Summary of the quantum data package.}
Starting from the deterministic regulator gauge package of \cite{continuumqft},
this paper has completed the exact finite-dimensional quantization step.  The
link variables were promoted from data-derived transporters to independent
$SU(3)$ variables, the full configuration space and product measure were
defined, and the partition function was proved finite for the free massive
scalar sector used here.  The gauge redundancy was solved globally by a rooted
maximal-tree gauge, with exact orbit uniqueness, exact Faddeev--Popov identity,
constant determinant, and a decoupled ghost sector.  The fermions integrate to
a finite determinant, the scalar sector is an exact Gaussian field with
stiffness operator $L_\eta+m_\phi^2I$, and the continuum action
$S_{\mathrm{YM}}^{\mathrm{cont}}$ (in the stationary formulation of
\cite{continuumqft}) is paired here with a canonical continuum-indexed class of
local gauge-invariant plaquette observables and its finite-regulator
approximation map.

\paragraph{Forward look to the thermodynamic theory.}
The $N\to\infty$ thermodynamic limit developed in the companion papers places
the fixed-$N$ partition functions constructed here into the stationary
mean-field setting: \cite{convergence} proves
\[
\pi_{N,1}\Rightarrow \pi_\infty,
\]
while \cite{continuumqft} identifies the continuum limit action
$S_{\mathrm{YM}}^{\mathrm{cont}}$ and associated thermodynamic observables on
$\pi_\infty$.  The paper supplies the exact finite-dimensional
quantization step on the deterministic regulator lattice, together with the
local Haar-coordinate analysis and the explicit gauge fixing needed to pass from
data package $\mathcal Q^{(N)}$ to its quantum ensemble.  In addition, it fixes
the continuum-indexed gauge-invariant observable class that is local on each
finite regulator, uniformly bounded on compact windows, nontrivial, and carried
in the thermodynamic presentation by the equilibrium product measures
\(\pi_\infty^{\otimes k}\) with the transferred logarithmic Sobolev control.

\bibliographystyle{plain}
\bibliography{references}

\end{document}
